\documentclass[10pt,letterpaper]{article}
\usepackage[T1]{fontenc}
\usepackage[utf8]{inputenc}
\usepackage{lmodern}
\usepackage[letterpaper,left=.91in,right=.91in,top=.82in,bottom=.81in,headheight=14pt]{geometry}
\usepackage{amsmath,amssymb,amsthm,mathtools,booktabs,array,microtype}
\usepackage{enumitem,fancyhdr,xcolor,titlesec}
\definecolor{ink}{HTML}{183D47}
\definecolor{accent}{HTML}{17635D}
\usepackage[colorlinks=true,linkcolor=accent,urlcolor=accent,citecolor=accent,bookmarksnumbered=true]{hyperref}
\usepackage{bookmark}
\hypersetup{pdftitle={Prime counting in intuitionistic primitive recursive arithmetic},pdfsubject={Exact formulas, sharp inequalities, asymptotic estimates and Lean formalization}}
\pagestyle{fancy}\fancyhf{}\fancyhead[L]{\small\color{ink}Prime counting in iPRA}
\fancyhead[R]{\small\color{ink}\thepage}\renewcommand{\headrulewidth}{.3pt}
\setlength{\emergencystretch}{2em}
\setlength{\parindent}{1.15em}
\setlength{\parskip}{2pt}
\allowdisplaybreaks[2]
\setcounter{tocdepth}{1}
\titleformat{\section}{\Large\bfseries\color{ink}}{\thesection}{.65em}{}
\titleformat{\subsection}{\large\bfseries\color{ink}}{\thesubsection}{.6em}{}
\titlespacing*{\section}{0pt}{12pt plus 3pt}{6pt}
\titlespacing*{\subsection}{0pt}{8pt plus 2pt}{4pt}
\setlist{topsep=4pt,itemsep=3pt,parsep=0pt}
\newtheorem{theorem}{Theorem}[section]
\newtheorem{lemma}[theorem]{Lemma}
\newtheorem{proposition}[theorem]{Proposition}
\newtheorem{corollary}[theorem]{Corollary}
\theoremstyle{remark}\newtheorem{remark}[theorem]{Remark}
\numberwithin{equation}{section}
\begin{document}
\thispagestyle{empty}
{\noindent\color{accent}\sffamily\small CONSTRUCTIVE NUMBER THEORY \hfill NOVEMBER 22, 2026\par}
\vspace{9pt}
{\noindent\fontsize{25}{29}\selectfont\bfseries\color{ink}Prime counting in intuitionistic\\primitive recursive arithmetic\par}
\vspace{8pt}
\noindent We prove exact sieve and inversion formulas for $\pi(x)$, determine the global maximum of $\pi(x)\log(x)/x$, and certify the exact cutoff $2418$ in Ramanujan's inequality for the logarithmic integral. Each assertion about an integral or a limit is read as an estimate for rational approximations with a primitive-recursive accuracy schedule. A pinned ProveIt derivation supplies the prime number theorem. A separate quantitative human proof supplies the stronger input for the asymptotic expansions of $\pi(x)$ and $p_m$. We recompute fourteen rows of the prime-count tables, and give counterexamples to numerical cutoffs in the cited material. The final section identifies the checked Lean theorems and the remaining object-level derivations.
\smallskip
\begingroup\small\setlength{\parskip}{0pt}\tableofcontents\endgroup
\newpage
\section{What a prime count means in iPRA}
\label{sec:found}
Our theory is first-order intuitionistic arithmetic over the natural numbers with symbols for the primitive-recursive functions, their defining equations, and induction for quantifier-free formulas. A rational number means a signed numerator and a positive denominator. Two rationals are equal when the cross products agree. Division of natural numbers used below always has an explicit positive divisor on its mathematical domain. A default value outside that domain makes each algorithm total.

For a decidable test $B$ write $[B]\in\{0,1\}$ for its truth value. A bounded universal claim for $i<t$ in an induction invariant is a product of $t$ such bits. Existence over the same interval is a bounded search, recording the first successful index and a success bit. These folds and bounded sums, products, minima and maxima are defined by primitive recursion. For a finite array of length $r$ with entries $<b$, $b\ge2$, use $\sum_{i<r}a_i b^i$ with bound $b^r$; reading a digit uses quotient and remainder. A scan with decreasing argument is implemented with a prescribed number of iterations or by building the table up to the original input. All finite sums over primes are thus sums over integers with a prime bit. Parameters such as a precision or an input ceiling are fixed before induction on an array. In particular a statement needed for all smaller inputs can be kept as a single prefix-test bit.

Here are the conventions for analysis. An assertion about a named real number is an assertion about a specified sequence of rational enclosures, with a primitive-recursive width schedule. Equality means agreement within every requested tolerance. Non-strict inequalities permit arbitrarily small rational slack. A strict conclusion supplies a positive rational separation, possibly depending on the inputs. A limit with error $1/k$ has a primitive-recursive threshold $T(k)$ and a proof valid for every input beyond it. These meanings are first-order schemes or universal closures of statements about numbers. For discontinuous counts we use integer or rational cuts, roots decided by powering, or an integer part together with boundary comparisons supplied in the input. An arbitrary fast Cauchy name alone need not decide equality to a prime. On a region with a supplied distance from the jumps the same step function is evaluated locally by a rational approximation of sufficient precision.

\subsection{Division, factorization and bounded prime selection}

For $n\ge0$ set
\[
 b(n)=[n\ge2]\prod_{d=2}^{n-1}[n\bmod d\ne0],
 \qquad \pi(n)=\sum_{j=1}^n b(j).
\]
Empty products are one and empty sums zero. This is the count including the upper endpoint. For $x=A/B\ge0$, $B>0$, put $\pi(x)=\pi(\lfloor A/B\rfloor)$. Extend it by zero on negative rationals. Write $p_1,p_2,\ldots$ in increasing order, with indices starting at one.

For $m>1$, take the least divisor $d$ in $[2,m]$ by bounded search. If $d$ had a proper divisor above one it would contradict minimality. The step $m\mapsto m/d$ strictly decreases the positive quotient whenever it exceeds one. In at most $m$ steps it reaches one, so the recorded factors multiply to $m$. They are nondecreasing because any smaller prime in the quotient would already have divided the preceding dividend. The product of the factors found so far times the current quotient equals the original input; this equality is the induction invariant.

For uniqueness we use Euclid's lemma. The Euclidean algorithm on $p,a$ tracks a Bézout representation of both current remainders. Start with coefficient pairs $(1,0),(0,1)$. A quotient step with quotient $q$ replaces the pairs $U,V$ by $V,U-qV$. Remainders decrease after the first division, so $a+p+2$ iterations, stopping when the second remainder is zero, suffice. Common divisors before and after each step agree by the remainder identity. At termination the last positive remainder divides both arguments. For prime $p$ with $p\nmid a$ it must be one, yielding $up+va=1$. Multiplying by $c$ shows $p\mid ac$ implies $p\mid c$ in that case. The coefficient magnitudes are bounded by $(a+p+2)^{a+p+2}$ because each step grows the previous maximum by at most $1+q\le a+p+2$. Hence the whole argument is a bounded computation with signed integers. Euclid's lemma applied successively to a finite product locates a matching prime on the other side of an equality of prime products. Using sorted lists, the least prime on either side is the same; cancel and apply length induction. This proves existence and uniqueness of the sorted factorization, including the empty factorization of one.

\begin{theorem}[Rank and inverse rank]
The function $\pi$ has jumps of size $b(n)$, is nondecreasing, and satisfies
\[
 \pi(p_k)=k\quad(k\ge1),\qquad
 p_{\pi(n)}\le n<p_{\pi(n)+1}\quad(n\ge2),\qquad
 p_{\pi(n)}=n\ \Longleftrightarrow\ b(n)=1\quad(n\ge2).
\]
In particular its integer range is all of $\mathbb N$.
\end{theorem}
\begin{proof}
The sum gives $\pi(0)=0$ and $\pi(n+1)=\pi(n)+b(n+1)$. Each increment is either zero or one, and $\pi(n)\le n$. To put a uniform bound on inverse search, take $E(0)=1$, $E(k+1)=E(k)!+1$. The least divisor $q>1$ of $E(k)!+1$ is prime and exceeds $E(k)$ since a divisor between two and $E(k)$ would divide both $E(k)!$ and its successor. Thus $\pi(E(k))\ge k$ by quantifier-free induction. Define $p_k$ for $k\ge1$ to be the least $j\le E(k)$ with $\pi(j)\ge k$. Since every jump is at most one, $\pi(j-1)=k-1$, $b(j)=1$ and $\pi(j)=k$. Monotonicity now gives $p_k\le n$ exactly when $k\le\pi(n)$, also if $n>E(k)$. For $n\ge2$ the count is positive, so applying this equivalence to its rank and the next rank yields the inequalities. A flagged point is the first occurrence of its positive rank, and every first occurrence is flagged. Zero occurs at zero and one.
\end{proof}

It follows also that $\pi$ is constant on $[p_k,p_{k+1})$, has left value $k-1$ and right value $k$ at $p_k$, and satisfies $0\le\pi(v)-\pi(u)\le v-u$ for integers $0\le u\le v$. Its initial values for $n=1,\ldots,16$ are $0,1,2,2,3,3,4,4,4,4,5,5,6,6,6,6$, from the primes $2,3,5,7,11,13$. Arbitrarily long constant stretches also have PR witnesses. For $s\ge1$, each integer $(s+2)!+j$ for $2\le j\le s+2$ has the proper divisor $j$.

\subsection{Congruence classes}
For $a\ge1$ and $0\le r<a$, the modular prime-counting function of MathWorld is
\[
 \pi_{a,r}(N)=\sum_{m=1}^N b(m)[m\bmod a=r].
\]
The division algorithm identifies this class with the integers $ak+r$ with $k\ge0$. Every prime contributes to exactly one class, so $\sum_{r<a}\pi_{a,r}(N)=\pi(N)$. For $g=\gcd(a,r)>1$, every class member is divisible by $g$, and a prime in the class must equal $g$. Thus $\pi_{a,r}(N)=b(g)[g\le N][g\bmod a=r]$. More generally
\[
 \sum_{\substack{0\le r<a\\\gcd(a,r)=1}}\pi_{a,r}(N)
  =\pi(N)-\sum_{p\le N}[p\mid a],
\]
where the last sum is over primes. Indeed $\gcd(a,m)=\gcd(a,m\bmod a)$, and the only primes having a common divisor above one with $a$ are its prime divisors. This includes $a=1$ with the unique residue zero.

\section{Counting the survivors of a finite sieve}
\label{sec:survivors}

For integers $N,a\ge0$, let $\Phi(N,a)$ count the integers in $[1,N]$ divisible by none of the first $a$ primes $p_1,p_2,\ldots,p_a$. Thus $\Phi(0,a)=0$ and $\Phi(N,0)=N$. For $r\ge1$ write $R_r(N)=\max\{t\le N:t^r\le N\}$. It satisfies
\[
 t^r\le N\quad\Longleftrightarrow\quad t\le R_r(N).
\]
Every endpoint in this section can therefore be calculated by integer comparisons. Write $N/d$ inside $\Phi$ or $\pi$ for the integer quotient $\lfloor N/d\rfloor$ when $d>0$. Prime indices start at one.

\subsection{The marking algorithm and Legendre's formula}

If $m\ge2$ is composite, its least factor $q>1$ is prime and $q^2\le m$. Indeed, a proper divisor of $q$ would be a smaller factor of $m$, and in $m=qu$ the complementary factor satisfies $u\ge q$. Consequently the following sieve produces exactly the primes through $N$. Set a flag at each $m$ from $2$ to $N$. In increasing order $j=2,\ldots,R_2(N)$, whenever $j$ is flagged, clear the flags at $j^2,j^2+j,j^2+2j,\ldots$ through $N$. Every cleared integer is composite. A composite $j$ has already been cleared when its least prime factor was processed, since that prime remains flagged and $j$ is one of its indicated multiples. The same argument clears every composite $m\le N$ by termination. The flags are a word of $N+1$ bits, and either loop uses at most $N$ iterations. This proves the correctness of the sieve of Eratosthenes in the required arithmetic.

\begin{proposition}[Legendre's identity and sieve recursion]
For $N\ge1$ and $b=\pi(R_2(N))$,
\[
 \pi(N)+1=\Phi(N,b)+b.
\]
For $a\ge1$ and $N\ge0$,
\[
 \Phi(N,a)+\Phi(N/p_a,a-1)=\Phi(N,a-1).
\]
\end{proposition}
\begin{proof}
The survivors after the first $b$ primes are $1$ and the primes in $(R_2(N),N]$. For the second identity, the integers lost at step $a$ are exactly $p_a u$ with $1\le u\le N/p_a$ surviving step $a-1$. For $i<a$, Euclid's lemma gives $p_i\mid p_a u$ exactly when $p_i\mid u$. Multiplication by $p_a$ and exact division are inverse on these two finite lists.
\end{proof}

For a list of \emph{any} distinct primes $q_1,\ldots,q_a$, let $F(N)$ be its survivor count. Code a subset by $s<2^a$ and let $d_s$ be the product of those $q_i$ whose bit is set, with $d_0=1$. Let $w(s)$ count the set bits. Then
\begin{equation}
 F(N)=\sum_{s<2^a}(-1)^{w(s)}\left\lfloor\frac{N}{d_s}\right\rfloor .
 \label{eq:sv-ie}
\end{equation}
For each $m\in[1,N]$ expand $\prod_{i=1}^a(1-[q_i\mid m])$. This product is its survival indicator. A subset in the expansion contributes $(-1)^{w(s)}[d_s\mid m]$ by unique factorization. Summing first over $m$ gives \eqref{eq:sv-ie}, since $d_s,2d_s,\ldots,\lfloor N/d_s\rfloor d_s$ are exactly its positive multiples through $N$. The product expansion is an induction on prefix length, splitting the next bit into 0 and 1. This proves the binary-mask formula underlying Mapes' method. For the first $a$ primes it gives $\Phi(N,a)$, with strictly increasing indices inside each subset. Repeated indices enter a different count below.

Two useful exact consequences follow. Put $Q=\prod_i q_i$ and $D=\prod_i(q_i-1)$. Since each $d_s$ divides $Q$, summing the change of the floors gives
\[
 F(N+Q)-F(N)=\sum_{s<2^a}(-1)^{w(s)}Q/d_s
 =Q\prod_{i=1}^a(1-1/q_i)=D.
\]
In particular $F(tQ+u)=tD+F(u)$ whenever $0\le u<Q$. For $a\ge1$ let $\delta=D/Q$. Replacing each floor in \eqref{eq:sv-ie} by its rational argument changes the contribution of an even mask downwards by less than one, and that of an odd mask upwards by less than one. Toggling the first bit pairs the masks of opposite parity. Hence
\begin{equation}
 |F(N)-N\delta|\le 2^{a-1}.
 \label{eq:sv-floor-bound}
\end{equation}
Rational calculations here and below may be multiplied by their displayed positive denominators.

\subsection{Meissel and Lehmer by sorting factorizations}

Let $P_k(N,a)$ count the integers $m\in[1,N]$ whose factorization contains exactly $k$ primes, with multiplicity, all of index $>a$. Take the empty factorization of $1$. For $N\ge1$, $P_0(N,a)=1$. If $0\le a\le\pi(N)$, then $P_1(N,a)=\pi(N)-a$. In general the latter formula uses truncated subtraction. Each $m$ has exactly one sorted factor list. Its length satisfies $2^k\le m$. Thus, for $N\ge1$ and $L=\max\{j\le N:2^j\le N\}$,
\[
 \Phi(N,a)=\sum_{k=0}^L P_k(N,a).
\]
For $N=0$ all these counts, including $P_0$, are zero.

For $k\ge0$ and $N\ge1$ splitting off the first prime gives the more general recurrence
\begin{equation}
 P_{k+1}(N,a)=\sum_{a<i\le \pi(R_{k+1}(N))}
       P_k(N/p_i,i-1).
 \label{eq:sv-factor}
\end{equation}
In any sorted product $p_i q_1\cdots q_k$ we have $p_i^{k+1}\le N$ and each $q_j\ge p_i$. Conversely, inserting $p_i$ at the front of any list counted on the right reconstructs an admissible list of length $k+1$. The unrestricted repetitions of $p_i$ explain the index $i-1$ on the right. The sets for different $i$ are disjoint because least prime factors are unique. This proof covers $k=0$ and perfect powers.

Let $b=\pi(R_2(N))$, $c=\pi(R_3(N))$ and $b_i=\pi(R_2(\lfloor N/p_i\rfloor))$. Applying \eqref{eq:sv-factor} once and twice gives
\begin{align}
 P_2(N,a)&=\sum_{a<i\le b}\big(\pi(N/p_i)-i+1\big),\label{eq:sv-p2}\\
 P_3(N,a)&=\sum_{a<i\le c}\ \sum_{i\le j\le b_i}
     \big(\pi(N/(p_i p_j))-j+1\big).\label{eq:sv-p3}
\end{align}
For example, at a fixed pair $i,j$ in the latter sum, the last factor has index in $[j,\pi(N/(p_i p_j))]$. This interval has the stated positive length by $p_ip_j^2\le N$. The integer quotient in two successive divisions is the quotient by their product, as follows by the division algorithm. Empty sums equal zero.

\begin{proposition}[Meissel's cube-root cutoff]
Take $N\ge1$, let $R_3(N)\le y\le R_2(N)$ and put $a=\pi(y)$, $h=b-a$. Then
\[
 \pi(N)=\Phi(N,a)+a(h+1)+\frac{h(h-1)}2-1
               -\sum_{i=a+1}^b\pi(N/p_i).
\]
In particular $a=c$ gives Meissel's formula.
\end{proposition}
\begin{proof}
Every surviving prime is at least $y+1$, whose cube exceeds $N$. Thus $P_k=0$ for $k\ge3$. The partition gives $\pi(N)=\Phi(N,a)+a-1-P_2(N,a)$. In \eqref{eq:sv-p2}, the sum of $i-1$ over $i=a+1,\ldots,b$ is $ha+h(h-1)/2$. Substitution proves the result.
\end{proof}

\begin{proposition}[Lehmer's fourth-root cutoff]
For $N\ge1$, set $a=\pi(R_4(N))$. With $b,c,b_i$ as above,
\begin{align*}
 \pi(N)={}&\Phi(N,a)+\frac{(b+a-2)(b-a+1)}2
       -\sum_{a<i\le b}\pi(N/p_i)\\
       &-\sum_{a<i\le c}\ \sum_{i\le j\le b_i}
                \big(\pi(N/(p_ip_j))-j+1\big).
\end{align*}
The constant term uses signed integer arithmetic, also when $N$ is small.
\end{proposition}
\begin{proof}
Now each surviving prime exceeds $R_4(N)$, so four cannot occur in one surviving product. Isolating $P_1$ gives $\pi(N)=\Phi(N,a)+a-1-P_2-P_3$. Substitute \eqref{eq:sv-p2} and \eqref{eq:sv-p3} and use
\[
 a-1+\sum_{a<i\le b}(i-1)=\frac{(b+a-2)(b-a+1)}2.
\]
This follows from the consecutive-integer sum or by induction on $b-a$. When $b=a$ the coefficient is $a-1$.
\end{proof}

For instance at $N=100$ the fourth-root cutoff is $3$ and $a=2$. The $33$ survivors have no factor $2$ or $3$, and their only composite products are $25,35,49,55,65,77,85,91,95$. Thus $P_2=9$, $P_3=0$, and $\pi(100)=33+2-1-9=25$. At $N=125$ the triple $5^3$ enters exactly at the inclusive boundary.

All the algorithms can be implemented by bounded folds. For fixed input ceiling $N$ the recursion for $\Phi$ fills rows $0,\ldots,a$ at columns $0,\ldots,N$. Each entry is at most $N$. Encode a row in base $N+2$; reading the two entries of the preceding row determines the next value. For \eqref{eq:sv-factor}, alternatively compute $P_k$ by scanning $m\le N$, repeatedly dividing by its least prime factor for at most $N$ steps and counting factors. In the identities, evaluation of $\pi$ at a quotient $N/p_i$ involves a smaller input. Using a table of previous values makes induction on $N$ sufficient if one wants the recursive evaluator. For $N\ge2$ all root indices needed in the cube- and fourth-root formulas also come from smaller arguments. The cases $N=0,1$ can be initialized directly. The combinatorial formulas and their proofs underlie the treatment by Lagarias, Miller and Odlyzko, \emph{Computing $\pi(x)$: the Meissel--Lehmer method}, Math. Comp. 44 (1985), 537--560.

\subsection{A factorial expression}

MathWorld's equation (7) admits a short proof that explains its endpoint. Put $r(j)=(j-2)!\bmod j$ for $j\ge3$, the unique residue between $0$ and $j-1$. For every prime $p\ge3$, multiplication by a nonzero residue permutes $1,\ldots,p-1$. Indeed, equality of products modulo $p$ forces equality of the factors modulo $p$ by Euclid's lemma, and an injective self-map of a finite interval is surjective by counting. Each residue therefore has a unique inverse, found by scanning this interval. Inversion pairs all residues except its fixed points. If $u^2\equiv1\pmod p$ then $p\mid(u-1)(u+1)$, giving $u=1$ or $u=p-1$. The paired residues each contribute $1$ and the fixed ones contribute $-1$. This proves Wilson's congruence $(p-1)!\equiv-1\pmod p$. Dividing off the factor $p-1\equiv-1$ shows $r(p)=1$.

Suppose $j\ge4$ is composite. Write $j=uv$ with $2\le u\le v$. If $u<v$ then $v\le j/2\le j-2$ and the factorial contains the distinct factors $u,v$. If $u=v\ge3$, it contains the distinct factors $u,2u$, since $2u\le u^2-2$. In both cases $r(j)=0$. The remaining case has $r(4)=2$. Thus, for $N\ge4$,
\[
 \pi(N)=-1+\sum_{j=3}^N r(j)
        =-1+\sum_{j=3}^N\left((j-2)!-j\left\lfloor
                       \frac{(j-2)!}{j}\right\rfloor\right).
\]
The sum counts the primes $3,\ldots,N$ once and has the additional contribution $2$ at $4$. This accounts for both the constant and the lower endpoint.

\section{Every integral value of the ratio}
\label{sec:golomb}

MathWorld tabulates positive solutions of $N=k\pi(N)$ for $N\ge2$. We give a primitive-recursive bound for \emph{every} solution, a witness for every $k\ge2$, and complete solution lists for $2\le k\le7$.

\begin{lemma}[An effective density bound]
For every $k\ge1$ let $y=2^{4k}$, $a=\pi(y)$, and
\[
 B(k)=2k(a+2^a)+1 .
\]
For all $N\ge B(k)$ one has $k\pi(N)<N$.
\end{lemma}
\begin{proof}
For the primes through $y$ set $\delta=\prod_{p\le y}(1-1/p)$. Unique factorization embeds each $m\le y$, $m\ge1$, as a distinct exponent tuple in the expansion of $\prod_{p\le y}\sum_{e=0}^y p^{-e}$. The bound $e\le y$ follows from $2^e\le m$ for nonzero exponents. All terms are positive, so
\[
 \delta\sum_{m=1}^y\frac1m
 \le\prod_{p\le y}\left((1-1/p)\sum_{e=0}^y p^{-e}\right)
 =\prod_{p\le y}(1-p^{-(y+1)})\le1 .
\]
For $0\le j<4k$, the $2^j$ terms from $m=2^j+1$ through $2^{j+1}$ add to at least $1/2$. Including $m=1$ gives $\sum_{m\le y}1/m\ge1+2k$. Hence $k\delta\le1/2$. Every prime through $N$ is either among the $a$ sieving primes or is counted by $\Phi(N,a)$. From \eqref{eq:sv-floor-bound},
\[
 k\pi(N)\le k(a+2^a)+kN\delta\le k(a+2^a)+N/2<N
\]
for $N\ge B(k)$. The tuple expansion uses at most $(y+1)^a$ terms and all rational equalities are finite.
\end{proof}

\begin{proposition}[Golomb's theorem with bounded witnesses]
For $k\ge2$ the equation $N=k\pi(N)$ has a solution with $2\le N<B(k)$. The least solution, the largest solution, and the total number of solutions with $N\ge2$ are primitive-recursive functions of $k$ on this domain.
\end{proposition}
\begin{proof}
Set $f(N)=N-k\pi(N)$ as a signed integer. Then $f(2)=2-k\le0$ and $f(B(k))>0$. Each increment $f(j+1)-f(j)$ is either $1$ or $1-k$. If $f(2)=0$, take $N=2$. Otherwise scan for the first $j\in[3,B(k)]$ with $f(j)\ge0$. Since $f(j-1)\le-1$, the upper bound of $1$ on a single increment gives $f(j)=0$. Necessarily $j<B(k)$. A scan of the entire interval finds all solutions, by the lemma. This argument supplies a bounded version of S.~W.~Golomb's \emph{On the ratio of $N$ to $\pi(N)$}, Amer. Math. Monthly 69 (1962), 36--37.
\end{proof}

For practical exact enumeration choose an arbitrary $a\ge1$ with $k\delta<1$ for the first $a$ primes. The sharper error in \eqref{eq:sv-floor-bound} proves that all $N$ satisfying the equation lie strictly below
\[
 T(k,a)=1+\left\lfloor\frac{k(a+2^{a-1})}{1-k\delta}\right\rfloor .
\]
Indeed $\pi(N)\le a+N\delta+2^{a-1}$ implies $N\le k(a+2^{a-1})/(1-k\delta)$ at a solution. The following table, including the exponents and cutoffs, is checked by \texttt{golomb\_check.py}. It constructs the sieve through the largest cutoff, accumulates flags to obtain $\pi(N)$ and records each exact equality. Thus the finite check and the proved cutoff together establish each complete list.
\begin{center}
\small
\begin{tabular}{r r r l}
$k$&$a$&$T(k,a)$&All solutions $N\ge2$\\ \hline
2&2&25&2, 4, 6, 8\\
3&3&106&27, 30, 33\\
4&5&498&96, 100, 120\\
5&7&3\,646&330, 335, 340, 350, 355, 360\\
6&10&59\,867&1008, 1080, 1092, 1116, 1122, 1128, 1134\\
7&15&3\,948\,665&3059, 3066, 3073, 3080, 3087, 3094
\end{tabular}
\end{center}

\section{Prime powers, midpoint values and finite inversion}
\label{sec:powers}
At a prime square the prime-power count has a jump even though the prime count is locally constant. Taking an integer part too soon in the inversion formula can therefore give the wrong answer. We encode each endpoint by an integer comparison.

Fix a nonnegative rational $x=A/B$ with $B>0$. For $0\le h\le1$ rational put
\[
 e_h(u;x)=[Bu<A]+h[Bu=A].
\]
Brackets are the $0,1$ values of decidable tests. For a positive exponent $d$ set
\[
 C_h(d;x)=\sum_{p\le\lfloor x\rfloor} e_h(p^d;x),\qquad
 J_h(d;x)=\sum_{k=1}^{L}\frac{C_h(dk;x)}{k},
\]
where $p$ always ranges over primes and $L$ is the largest positive $j\le A$ with $B2^j\le A$, defaulting to zero. A term with $dk>L$ vanishes. These formulas evaluate counts at $x^{1/d}$ even when the root is irrational. The subscripts $1$ and $1/2$ give respectively the right-continuous and midpoint conventions. In the notation of the source pages, $C_{1/2}(1;x)=\pi_0(x)$ and $J_{1/2}(1;x)=\Pi_0(x)=f_0(x)$.

\begin{theorem}[Inversion at every endpoint]
For every $d\ge1$,
\[
 C_h(d;x)=\sum_{k=1}^L \frac{\mu(k)}{k} J_h(dk;x).
\]
For $x=p^a$, the jump of $J_1(1;x)$ has size $1/a$, and its midpoint value is the right value minus $1/(2a)$. For every other positive rational endpoint the left and right values agree, except at the prime powers just specified.
\end{theorem}
\begin{proof}
Use $\mu(1)=1$, $\mu(k)=0$ when some prime square divides $k$, and otherwise $\mu(k)=(-1)^r$ for $r$ distinct prime factors. This is a bounded computation, by trial division through $k$. The squarefree divisors of $j$ correspond to subsets of its distinct prime factors. For $j>1$ toggle its least prime in the subset. This fixed-point-free pairing gives
$\sum_{k\mid j}\mu(k)=[j=1]$; for $j=1$ the empty subset contributes one. The bijection uses uniqueness of factorization. One bounded way to obtain it is to divide by the least divisor above one until the quotient is one (at most $j$ steps). Each least divisor is prime. Euclid's lemma then identifies the first factor in any sorted prime factorization and cancellation gives uniqueness by induction on the length.
Expanding the right side of the displayed inversion and grouping pairs $(k,l)$ by $j=kl$, the coefficient of $C_h(dj;x)$ is $j^{-1}\sum_{k\mid j}\mu(k)$ whenever $j\le L$. All other summands vanish because $d j>L$. This leaves exactly $C_h(d;x)$, including when $L=0$ as both sides then vanish. For $h=1/2$, multiplication by $2(L!)^2$ clears every denominator in the expanded identity.
Unique factorization makes the map $(p,a)\mapsto p^a$ injective for prime $p$ and $a\ge1$. The change between the strict and weak comparisons therefore has exactly the asserted size. On both sides of the rational cut, any positive displacement smaller than $1/(2B)$ crosses no other integer. Thus the arithmetic midpoint equals the usual symmetric limit. The same observation makes $\pi$ locally constant, with derivative zero, at every rational nonprime point.
\end{proof}

For example $J_1(1;8)=4+1/2+1/3=29/6$ and $J_{1/2}(1;8)=14/3$. At the square-root cut, $J_{1/2}(2;8)=1$; the prime $2$ has full weight because $4<8$. At the cube-root cut $J_{1/2}(3;8)=1/2$. Inversion gives $14/3-1/2-1/6=4$. At $x=9$ the square-root cut gives prime $3$ half weight. Both cases follow from the same comparison.

\paragraph{The checked inversion theorem.}
The companion \texttt{Inversion.lean} defines $\mu$ by $\mu(1)=1$ and the negative sum of the proper-divisor values at $n>1$. The recurrence singles out the same function as the prime-factor definition above: subtract its proper-divisor sum from $\sum_{k\mid n}\mu(k)=[n=1]$, then use uniqueness of the recurrence (bounded strong induction). Its checked theorem \texttt{midpoint\_inversion} uses $L>0$ and arbitrary Boolean prime indicators. Put $D=L!$ and $w_j=(D/j)\,2C_{1/2}(dj;x)$ for $1\le j\le L$. The theorem regroups
\[
 \sum_{k=1}^L\mu(k)\sum_{\substack{j\le L\\k\mid j}}w_j=w_1
 \qquad (j\ge1).
\]
Here $D/j$ is exact integer division. In the inner sum write $j=kr$; the unique $r$ runs from $1$ to $\lfloor L/k\rfloor$. Thus this sum is $(2D/k)J_{1/2}(dk;x)$ for the cutoff $L$ specified above, while $w_1=2D C_{1/2}(d;x)$. This gives the rational reading of the checked integer equality. At $L=0$ the natural-bound sums vanish directly.

\section{Chebyshev functions and the two partial summations}
\label{sec:weights}
There is an integer version of every basic weighted counting identity. Write
\[
 U(n)=\prod_{p\le n}p,\quad V(n)=\operatorname{lcm}(1,\ldots,n),\quad
 D(m)=\begin{cases}p&m=p^a\ (a\ge1),\\1&\text{otherwise}.\end{cases}
\]
Empty products and the empty least common multiple are one. Write $r_d(n)$ for the largest integer $t\le n$ with $t^d\le n$, for $d\ge1$. Taking the multiplicity of any prime $p\le n$ shows
\[
 V(n)=\prod_{m=2}^nD(m)=\prod_{d=1}^{L}U(r_d(n)),\qquad L=\lfloor\log_2 n\rfloor\quad(n\ge1).
\]
Indeed, the largest exponent of $p$ among $1,\ldots,n$ is precisely the number of positive $a$ with $p^a\le n$. For $p>n$ every multiplicity is zero. Thus for $\Lambda(m)=\log D(m)$, $\vartheta(x)=\log U(\lfloor x\rfloor)$ and $\psi(x)=\log V(\lfloor x\rfloor)$ we obtain the source identities
\[
 \psi(x)=\sum_{m=2}^{\lfloor x\rfloor}\Lambda(m)
       =\sum_{d=1}^{L}\vartheta(x^{1/d}),\qquad
 \vartheta(x)=\sum_{p\le x}\log p.
\]
For $m\ge2$, $\Lambda(m)/\log m$ is $1/a$ at $p^a$ and zero elsewhere. Consequently its sum through $x$, with the boundary term weighted by $1/2$, is exactly $J_{1/2}(1;x)$. This last identity can be encoded using just rational numbers.

Here and below a logarithm of $v\ge1$ is represented by upper and lower rectangular sums of $1/t$ from $1$ to $v$. At rational $v$, dividing into $M$ intervals gives a gap at most $(v-1)/M$ by telescoping the endpoint differences. For general positive rational arguments use $\log v=-\log(1/v)$ below one. Splitting $[1,ab]$ at $a$ and scaling $t=au$ in the second interval gives $\log(ab)=\log a+\log b$ for $a,b\ge1$. The argument uses finite partitions, refined by their common endpoints. The same identity for positive rationals follows by clearing inverses. Also for $1\le a\le b$,
\[
 (b-a)/b\ \le\ \log b-\log a\ \le\ (b-a)/a.
\]
These rational enclosures give equality at every requested precision. For a sum of $K$ logarithms, ask each enclosure for width $2^{-j}/(K+1)$ to obtain overall width at most $2^{-j}$. Integer factors in front of a sum multiply this precision cost by their absolute values.

\begin{theorem}[Partial summation]
At every rational $x\ge2$ one has
\begin{align*}
 \vartheta(x)&=\pi(x)\log x-\int_2^x \pi(t)\frac{dt}{t},\\
 \pi(x)&=\frac{\vartheta(x)}{\log x}
        +\int_2^x\frac{\vartheta(t)}{t\log^2t}\,dt .
\end{align*}
\end{theorem}
\begin{proof}
Include all integers strictly between $2$ and $x$ as cut points, together with $2,x$, without repetitions. For any nondecreasing list of endpoints $u_i$, prefix sums $S_i$ of weights, and selected values $f_i$, direct cancellation gives the discrete product rule
\[
 S_i f_i-S_{i-1}f_{i-1}=(S_i-S_{i-1})f_i+S_{i-1}(f_i-f_{i-1}).
\]
For the first identity take $S_i=\pi(u_i)$, $f_i=\log u_i$. The changes in $S_i$ contribute $\log p$ at each prime, with the starting product $\pi(2)\log2$ supplying the first prime. The other terms equal the integrals over individual intervals since $\pi$ is constant in their interiors. For the second identity use $S_i=\vartheta(u_i)$, $f_i=1/\log u_i$; every jump product is one and $\vartheta(2)/\log2=1$. The complementary differences use
\[
 \frac{1}{\log a}-\frac{1}{\log b}
  =\frac{\log b-\log a}{\log a\log b}
  =\int_a^b\frac{dt}{t\log^2t}\qquad(2\le a\le b).
\]
For the last equality apply the logarithmic enclosure above on each mesh interval $[u,v]$. The decrement of $1/\log t$ is sandwiched between $(v-u)/(v\log^2 v)$ and $(v-u)/(u\log^2 u)$. At mesh width $h$, summing gives an uncertainty at most $2h$ by monotonicity and $\log2\ge1/2$. Thus the integral identities are quantitative assertions about rational sums. For explicit bookkeeping let $X=\lceil x\rceil$; the step heights are at most $X^2$ and there are at most $X$ intervals. Taking $h\le2^{-j-3}/(X+1)^3$ gives total quadrature error below $2^{-j-1}$. Computed reciprocal logarithms are safe because their denominators are bounded below by $1/2$.
\end{proof}

\section{The positive-real Mellin identity with a truncation bound}
\label{sec:mellin}
The Mellin identity for $\Pi_0$ on the Wikipedia page has an elementary proof on positive rational exponents $s>1$. This part lies in the region of absolute convergence. Its internal meaning is an estimate at finite cutoffs, with an explicit primitive-recursive accuracy schedule.

\begin{theorem}
Let $s=a/b$ with positive integers $a>b$. Write
\[
 Z_N=\sum_{n=1}^N n^{-s},\qquad
 T_N=\sum_{p^k\le N}\frac{p^{-ks}}{k},\qquad
 J(t)=J_{1/2}(1;t).
\]
For $K\ge1$ and $N=2^{bK}$ one has
\[
 \left|\log Z_N-s\int_1^N J(t)t^{-s-1}\,dt\right|
     \le (4b+1)2^{-K}.
\]
In particular, at the level of the associated names,
$\log\zeta(s)=s\int_0^\infty J(t)t^{-s-1}\,dt$, where $\zeta(s)$ is the name given by the $Z_N$. The integrand below $1$ is assigned zero.
\end{theorem}
\begin{proof}
For any $H>N$,
\[
 \sum_{N<n\le H}n^{-s}\le D_K:=2b\,2^{-K}.
\]
In a block $2^j\le n<2^{j+1}$ the sum is at most $2^{-j/b}$ since $s\ge1+1/b$. Group $b$ successive blocks starting at $j=bK$; their combined bound is $b2^{-K}$. Subsequent groups form a geometric series of ratio $1/2$. Every finite sum is bounded by twice the first group. In particular $Z_N$ is an effective Cauchy name. Rational powers in this proof have primitive-recursive enclosures: for $n^{a/b}$ bisect $[1,n^a]$, comparing the $b$th power of each rational midpoint with $n^a$; inversion on $[1,\infty)$ does not increase errors. For rational positive bases use the same method with rational endpoints widened to integers as needed.

Consider the finite product $E_N=\prod_{p\le N}(1-p^{-s})^{-1}$. To compare it with $Z_N$, first truncate each geometric factor after exponent $m\ge N$. On multiplying out, every $n\le N$ appears exactly once by unique factorization; each other monomial equals $n^{-s}$ for a distinct $n>N$. The resulting difference is between zero and $D_K$, even for the largest product of prime powers in this rectangular expansion. Passing from finite factors to $E_N$ is effective. Each full factor is at most $2$, and replacing it by its truncation costs at most $2^{-m}$ since $p^{-s}\le1/2$. A telescoping product estimate gives total error at most $N 2^N 2^{-m}$. It follows that $0\le E_N-Z_N\le D_K$. Since both are at least one, the logarithm enclosure gives $0\le\log E_N-\log Z_N\le D_K$.

For $0\le q\le1/2$, summing a geometric progression inside $\int_0^q(1-u)^{-1}du$ gives
\[
 0\le -\log(1-q)-\sum_{k=1}^{m}\frac{q^k}{k}
       \le 2q^{m+1}.
\]
Indeed the geometric remainder under the integral is $u^m/(1-u)$, bounded by $2u^m$; integrating powers of $u$ follows by telescoping $(u+h)^k-u^k$ between its lower and upper rectangular bounds. With $q=p^{-s}$ and $m\ge N$, the terms of the finite sums omitted by the cut $p^k\le N$ correspond to distinct integers $>N$. Their sum is at most $D_K$. The discarded geometric remainders contribute at most $N2^{-m}$. Additivity of the logarithm therefore gives $0\le\log E_N-T_N\le D_K$. In particular $|\log Z_N-T_N|\le2D_K$.

On each interval between successive integers, $J(t)$ agrees in its interior with the right-continuous sum $\sum_{p^k\le t}1/k$. Integrating each step separately gives the exact finite formula
\[
 s\int_1^N J(t)t^{-s-1}\,dt=T_N-N^{-s}\sum_{p^k\le N}\frac1k.
\]
For the endpoint kernel identity, take $1\le u<v$ and write $u=\alpha^b$, $v=\beta^b$. After canceling $\beta-\alpha$ in $(u^{-s}-v^{-s})/(v-u)$, its numerator is the sum of $a$ products $\alpha^{-i}\beta^{-(a+1-i)}$, $1\le i\le a$, and its denominator the sum of $b$ products $\beta^{b-1-i}\alpha^i$, $0\le i<b$. Bounding each factor by its endpoint values traps the quotient between $s v^{-s-1}$ and $s u^{-s-1}$. Summing on a mesh proves $s\int_u^v t^{-s-1}dt=u^{-s}-v^{-s}$; at mesh width $h$ the upper and lower sums differ by at most $sh$. This also justifies the finite formula using only rational-power enclosures. Unique prime powers give $\sum_{p^k\le N}1/k\le N$, hence the correction is at most $N^{1-s}\le2^{-K}$. This proves the inequality. The same estimates at arbitrary upper limits $H\ge N$ (the tail after $H$ is bounded by $D_K$) identify the improper integral with the stated name. A primitive-recursive choice sufficient for error $2^{-j}$ in the displayed bound is $K=j+4b+2$.
\end{proof}

In this proof, the finite product over integers up to $N$ is a bounded fold. Even the rectangular expansion at cutoff $m$ has at most $(m+1)^N$ entries, all supported at integers at most $N^{mN}$. To read the inequalities with irrational factors entirely as interval arithmetic, fix an additional error $2^{-j}$. Enclose each factor and refine all finite sums and products by their usual rational interval rules. The fold lengths and a bound on their local amplification give a primitive-recursive precision schedule: for a circuit of $S$ addition, subtraction, multiplication and reciprocal gates, with absolute values bounded by $M\ge2$ and inverted quantities bounded below in modulus by $1/M$, round the inputs to width at most $2^{-j-2}(8M^2)^{-S}$. Each gate on that enclosure amplifies the maximum incoming width by at most $8M^2$, once incoming widths are at most $1/(2M)$. If needed shrink the initial width by a further $2M$. Here a bound $M$ can be propagated from the rational bases, the explicit positive lower bounds on divisors, and the finite gate list. Arbitrary precision followed by the rational inequality with an arbitrarily small slack is the convention for a non-strict comparison of names.

\subsection{The Mellin identity for complex rational parameters}
The integral identity above extends to $s=\sigma+i\tau$ with rational $\sigma>1$ and rational $\tau$. Define $v^{-s}$ by $v^{-\sigma}(\cos(\tau\log v)-i\sin(\tau\log v))$. Choose $b\ge1$ with $\sigma\ge1+1/b$. Define the branch of $\log\zeta(s)$ by the convergent name
\[
 T(s)=\sum_{p,k\ge1}\frac{p^{-ks}}{k},
\]
where $p$ ranges over primes. At the cutoff $p^k\le N=2^{bK}$ the tail in complex modulus is at most $2b\,2^{-K}$: different prime powers have different integer bases of summation in $\sum_{n>N}n^{-\sigma}$. For $H\ge N$, the finite step integral gives
\[
 s\int_1^H J(u)u^{-s-1}\,du
   =\sum_{p^k\le H}\frac{p^{-ks}}k-H^{-s}\sum_{p^k\le H}\frac1k.
\]
The boundary correction has modulus at most $H^{1-\sigma}$. The kernel identity follows by differentiating $u^{-s}$ on each cell. For a rational verification on $[1,H]$, the second derivative has modulus at most $|s|(|s|+1)$. Taylor's formula with the resulting quadratic error gives the cell identity by telescoping. Thus the ordinary improper integral equals $T(s)$ with an explicit truncation error at most $(2b+1)2^{-K}$.

To identify the branch, let $T^{(N)}$ include every exponent for each $p\le N$. At $|z|\le1/2$ the identity $\exp(\sum_{k\ge1}z^k/k)=1/(1-z)$ follows by differentiation along $uz$, $0\le u\le1$: the derivative of $(1-uz)\exp(\sum_{k\ge1}(uz)^k/k)$ is zero and its initial value is one. The geometric tails bound the differentiated polynomials uniformly. Consequently $\exp T^{(N)}=\prod_{p\le N}(1-p^{-s})^{-1}$. The product of truncated geometric sums contains each $n\le N$ once by factorization, and the sum of the absolute values of the extra monomials is at most $2b\,2^{-K}$, as in Section~\ref{sec:mellin}. Letting the geometric depth increase, the same inequality holds between this finite Euler product and $Z_N=\sum_{n\le N}n^{-s}$. Since $|T^{(N)}-T|\le2b\,2^{-K}$ and $|T^{(N)}|,|T|\le1+2b$ by the dyadic bound from $1$, the exponential amplifies this error by at most $3^{2+2b}$, by integrating its derivative along the connecting segment. This proves $\exp T=\zeta(s)$ and therefore identifies the branch without a logarithm choice or a zero test. Complex numbers here are pairs of rational enclosure names. Sine, cosine and exponential on $|z|\le U$ use their Taylor polynomials: for $m\ge2U$ successive absolute majorants beyond $m$ decrease at least geometrically, giving error at most $2U^{m+1}/(m+1)!$. The general length schedule at the end of Section~\ref{sec:li-cut} also suffices.


\section{The correction integral over the trivial zeros}
\label{sec:trivial}
A further analytic identity from the Wikipedia page can be isolated from its explicit formula. For rational $x>1$ put
\[
 C(x)=\int_x^\infty\frac{dt}{t(t^2-1)\log t}.
\]
For $0<z<1$ define $\operatorname{li}(z)=\int_0^z du/\log u$. This integral uses only the part of the logarithmic integral below its singularity. Write $q=x^{-2}$, $\ell=(x-1)/x$. Then
\[
 0\le C(x)+\sum_{m=1}^M\operatorname{li}(x^{-2m})
    \le \frac{q^{M+1}}{2(M+1)(1-q)\ell}.
\]
In particular $C(x)=-\sum_{m\ge1}\operatorname{li}(x^{-2m})$ has an effective geometric rate, also for rational inputs close to one.

\begin{proof}
For $t\ge x$, $\log t\ge\ell$ by the lower rectangle for $\int_1^x dt/t$. Expand $1/(t^2-1)=\sum_{m=1}^M t^{-2m}+t^{-2(M+1)}/(1-t^{-2})$. The integral of the remainder is bounded above by
\[
 \frac1{(1-q)\ell}\int_x^\infty t^{-2M-3}dt
   =\frac{q^{M+1}}{2(M+1)(1-q)\ell}.
\]
On a finite interval $[x,T]$ substitute $u=t^{-2m}$ in each term; since $\log u=-2m\log t$ the oriented integral is
\[
 \int_x^T\frac{t^{-2m-1}}{\log t}\,dt
   =\int_{x^{-2m}}^{T^{-2m}}\frac{du}{\log u}.
\]
For a direct quadrature check, on $[v,w]\subseteq[x,T]$ the positive integral of $-1/\log u$ on the reversed image interval is between
$(v^{-2m}-w^{-2m})/(2m\log w)$ and $(v^{-2m}-w^{-2m})/(2m\log v)$. The endpoint kernel bounds of the preceding proof place these within the lower and upper rectangles for $t^{-2m-1}/\log t$ on $[v,w]$. Their summed gap tends to zero with mesh width, bounded by the width times the initial integrand value. Sending $T$ to infinity leaves $-\operatorname{li}(x^{-2m})$; the missing interval $[0,T^{-2m}]$ costs at most $T^{-2m}/(2m\log T)$. The tail of $C$ after $T$ costs at most $1/(2T^2(1-q)\ell)$ by the same remainder estimate with $M=0$ at the lower limit $T$.

All accuracy choices can be made by arithmetic on the input. For example put $D=\lceil 1/(1-q)\rceil$, $F=\lceil 1/((1-q)\ell)\rceil$. Bernoulli's inequality gives $q^D\le1/2$ because $(1/q)^D\ge1+D(1-q)$. For error $2^{-j}$ the choices $M=D(j+F+2)$ and $T=\lceil x\rceil+2^{j+F+2}$ make both indicated truncation errors individually smaller than $2^{-j-1}$. On a compact interval the positive integrands in the $t$ variable decrease; thus the rectangular gap is at most the mesh width times the initial value, bounded rationally by replacing $\log x$ with $\ell$. This supplies a primitive-recursive quadrature for the remaining finite integrals.
\end{proof}

\paragraph{The smooth arctangent correction.}
The same Wikipedia section describes a different correction,
\[
 A(x)=\frac{1}{\log x}-\frac1{\pi_{\rm c}}\arctan\frac{\pi_{\rm c}}{\log x},
 \qquad x>1,
\]
where $\pi_{\rm c}=4\arctan 1$ is the circle constant. Its vanishing at infinity follows without any property of Riemann's $R$ function. Define $\arctan z=\int_0^z(1+t^2)^{-1}dt$ for $z\ge0$ and put $z=\pi_{\rm c}/\log x$. Since $1/2\le(1+t^2)^{-1}\le1$ on $[0,1]$, we have $2\le\pi_{\rm c}\le4$. Integrate
\[
 0\le z-\arctan z=\int_0^z\frac{t^2}{1+t^2}dt\le z^3/3,
 \qquad
 0\le z^3/3-(z-\arctan z)=\int_0^z\frac{t^4}{1+t^2}dt\le z^5/5.
\]
Division by $\pi_{\rm c}$ gives $0\le A(x)\le\pi_{\rm c}^2/(3\log^3 x)$ and shows that the difference between this upper bound and $A(x)$ is at most $\pi_{\rm c}^4/(5\log^5 x)$. For $x\ge2^r$ the first bound is at most $128/(3r^3)$ since $\log2\ge1/2$. Thus $x\ge2^{2^{j+6}}$ suffices for $A(x)\le2^{-j}$. The defining arctangent integrand is decreasing; its rectangular enclosure has gap at most the mesh width, and endpoint uncertainties cost at most their width. This also gives an explicit arithmetic interpretation of the assertion that the approximation $R(x)-A(x)$ differs from $R(x)$ by a term tending to zero.

\section{Growth without analytic assumptions}
\label{sec:dyadicgrowth}

There are two routes to asymptotic information. An elementary binomial estimate already gives the correct order of magnitude. The leading constant $1$ will come from a specific previously proved iPRA theorem.

\begin{proposition}[Integer Chebyshev bounds]
Let $b$ be the unique integer with $2^b\le n<2^{b+1}$, for $n\ge1$. Then
\[
 \lfloor n/2\rfloor\le (b+1)\pi(n),\qquad
 b\pi(n)\le13n\quad(b\ge2).
\]
In particular $\pi(n)$ has order $n/\log n$, and $\pi(n)/n$ tends to zero with the explicit modulus $n\ge2^{13k+1}$ for error $<1/k$.
\end{proposition}
\begin{proof}
Write $L_n=\mathrm{lcm}(1,\ldots,n)$. For each $p$, counting the multiples of each $p^j$ in a factorial gives
\[
 v_p\binom{2m}{m}
 =\sum_{j=1}^{2m}\left(\lfloor2m/p^j\rfloor-2\lfloor m/p^j\rfloor\right).
\]
Each summand is either zero or one and vanishes beyond $p^j>2m$. This proves $\binom{2m}{m}\mid L_{2m}$. Every $p^j$ in $(m,2m]$ contributes a one, which also proves $L_{2m}/L_m\mid\binom{2m}{m}$. The inequalities
$2^m\le\binom{2m}{m}\le2^{2m}$ follow by multiplying $(m+i)/i\ge2$ for $1\le i\le m$, and by summing the entire binomial row. Hence $2^{\lfloor n/2\rfloor}\le L_n\le n^{\pi(n)}$ for $n\ge2$, the latter bound using the largest power of each prime separately. This proves the first inequality.

The second divisibility gives $L_{2^j}\le2^{2^{j+1}-2}$ by induction starting at $j=0$. Monotonicity gives $L_n\le2^{4n}$. Put $h=\lfloor b/2\rfloor$ and let $c$ count the primes in $(2^h,n]$. Their product divides $L_n$, so $hc\le4n$. The other primes number at most $2^h$. For $b\ge2$, $b\le3h$ and $b2^h\le2^b$: the latter follows from $b\le2^{\lceil b/2\rceil}$ by induction in steps of two starting at 2 and 3. Thus $b\pi(n)\le n+12n$. Every power, quotient, product and valuation here has a bounded computation.
\end{proof}

\section{The prime number theorem with an arithmetic modulus}
\label{sec:pntgrowth}

For $H_j=\sum_{i=1}^j1/i$ set $F(n)=(n^2)!$, $A(n)=\sum_{n<i\le n^2} F(n)/i$ and $\lambda(n)=A(n)/F(n)$. Every displayed division in $A$ is exact. For $n\ge1$ the right and left rectangle sums for $1/t$ between $n$ and $n^2$ give
\begin{equation}\label{eq:growth-harm}
 0\le\log n-\lambda(n)\le 1/n-1/n^2\le1/n.
\end{equation}
Here and below logarithms, integrals and limits of the specified functions abbreviate rational estimates at arbitrary requested precision. For instance right rectangles of width at most $2^{-s}$ on a positive interval $[1,x]$ give $\log x$ within $2^{-s}$, since the differences of the upper and lower rectangles telescope. For rational $x\ge1$ use equally spaced rectangles, $\max(1,\lceil (x-1)2^s\rceil)$ of them. For $0<x<1$ use $-\log(1/x)$. On $[2,x]$ one can compute $\mathrm{Li}_2(x)=\int_2^xdt/\log t$ to precision $2^{-s}$ by a mesh of width $2^{-s-3}$ (total rectangle error at most twice the mesh) and node logarithms to error $2^{-s-5}/(x+1)$, using $\log2\ge1/2$ and interval division. These procedures and their bounds are primitive recursive on rational codes.

\paragraph{The imported arithmetic theorem.}
At revision \texttt{a9c0f1949b}, ProveIt contains a deeply embedded derivation
\href{https://github.com/VladimirReshetnikov/ProveIt/blob/a9c0f1949be61679ed2e219b97dc42ecb7ae91a9/Logic/PrimitiveRecursiveArithmetic/Derivations/Lean/IPRADerivations/AnalyticNumberTheory/FinalPNT.lean}{\texttt{FinalPNT.arith\_pnt}}. With $k\ge1$ and $n\ge N_*(k)$ its conclusion is the integer inequality
\begin{equation}\label{eq:growth-embedded}
 k\,|\pi(n)A(n)-nF(n)|\le nF(n).
\end{equation}
The theorem constructs a derivation under \texttt{PRAAx} from derivations of the two displayed hypotheses. Its prime indicator uses trial division through the least prime factor. The least divisor $d\ge2$ of $n\ge2$ equals $n$ precisely when no $2\le t<n$ divides $n$, so its count agrees term by term with ours. This use of an existing theorem is the only prime-distribution input for the rest of this section.

Here is a simple larger modulus that can be used without reproducing the contraction polynomials. Put
\[
 q=5k,\quad D=2^{20}q^3,\quad V_0=D+1600,\quad
 V_{i+1}=[10^8(D+V_i+2)]^8,\qquad
 N(k)=2^{V_{4D}+131137q^3}.
\]
To check the comparison with the source, its \texttt{Nk} is $2^{s_J+131137q^3}$, $J=\lfloor(4D-2^{20}q^2)/3\rfloor$, $s_0=D+1600$, $A_j=4D-3j$. One stage sets $P=4(\lfloor24D/A_j\rfloor+2)+2$ and $s'=2(s+P(c+1))+2$ with $c=\texttt{DecLin}+\texttt{DecCst}$ in \href{https://github.com/VladimirReshetnikov/ProveIt/blob/a9c0f1949be61679ed2e219b97dc42ecb7ae91a9/Logic/PrimitiveRecursiveArithmetic/Derivations/Lean/IPRADerivations/AnalyticNumberTheory/StepCondition.lean}{\texttt{StepCondition}}. For $j<J$, $1\le A_j\le4D$, $P\le96D+10$. Writing $W=10^8(D+s+2)$, all factors $D,A_j,P,P+s+1$ and numerical coefficients in the expanded sums are at most $W$. Those sums contain fewer than 20 monomials, each a product of at most five such factors (including its coefficient). Thus $c\le20W^5$ and $s'\le W^8$. Induction up to $J\le4D$ gives $N_*\le N$. The underlying proof obtains a finite Selberg symmetry identity by divisor convolution, bounds the prime-power Chebyshev sum on successive blocks, improves the relative error by $3/D$ per stage, and then converts the prime-power sum to $\pi$ by finite summation and change of harmonic mesh. The stage predicate tests all integers from its threshold to a fixed endpoint, so the induction acts on a Boolean table.

\begin{theorem}[Two forms of the prime number theorem]
With $T(k)=\max\{N(2k),2k\}$, $k\ge1$, one has
\[
 n\ge T(k)\quad\Longrightarrow\quad
 \left|\frac{\pi(n)\log n}{n}-1\right|\le1/k .
\]
Both $\pi(n)\sim n/\log n$ and $\pi(n)\sim\mathrm{Li}_2(n)$ follow; the second statement also holds with the principal-value normalization $\mathrm{li}$.
\end{theorem}
\begin{proof}
Divide \eqref{eq:growth-embedded} at precision $2k$ by $nF(n)$ and use \eqref{eq:growth-harm} and $\pi(n)\le n$. The total normalized error is at most $1/(2k)+1/n$.

Write $L=\log n$. The identity
\[
 \mathrm{Li}_2(n)=n/L-2/\log2+\int_2^n\frac{dt}{(\log t)^2}
\]
comes from differentiating $t/\log t$. Equivalently, telescope its increments on rational mesh intervals and pass to the prescribed precision. The derivatives through order two of $t/\log t$ on $[2,n]$ have bounds $10(n+1)$, so a mesh of width $2^{-s}/(20(n+1)^2)$ suffices for the derivative-integral comparison. Split the integral at $\sqrt n$ into pieces bounded by $4\sqrt n$ and $4n/L^2$. It follows that
\[
 |L\mathrm{Li}_2(n)/n-1|\le4L/n+4L/\sqrt n+4/L.
\]
For $n\ge2^{512 k^2}$ we have $L\ge256 k^2$. Since $\exp(L/2)\ge L^3/48$, the last display is at most $384/L^2+4/L<1/(8k)$. Explicit comparisons with $\sqrt n$ may use the integer-root partition or rational enclosures to any fixed precision. The conventional principal-value normalization adds the constant $\operatorname{li}(2)$, of absolute value less than four by the explicit rational endpoint enclosure of Section~\ref{sec:li-cut}. Thus the normalization changes the last bound by at most $4L/n$. With, for instance, $n\ge\max(T(8k),2^{512(2k)^2})$, the normalized quantities for $\pi$, $\mathrm{Li}_2$ and $\mathrm{li}$ are all within $1/(4k)$ of one; division by a number at least $3/4$ gives ratio error at most $2/(3k)$ in either logarithmic-integral normalization.
\end{proof}

\section{Higher-order estimates from the logarithmic-integral bound}
\label{sec:growthrefine}

The basic PNT allows errors too large to decide even the sign of $\pi(n)-n/\log n$. For this section we use one additional result, the quantitative estimate in the \href{https://github.com/VladimirReshetnikov/ProveIt/tree/a9c0f1949be61679ed2e219b97dc42ecb7ae91a9/Logic/PrimitiveRecursiveArithmetic/docs/proofs/analytic-number-theory/pnt_li_ipra_2026-11-15}{ProveIt logarithmic-integral packet}, particularly \texttt{tex/analytic.tex}, \texttt{theta\_bridge.tex} and \texttt{transfer.tex}. This is a cited human proof with an arithmetic precision analysis. Its accompanying Lean source gives a specification. The previous section used the separate object-level theorem \texttt{FinalPNT.arith\_pnt}.

In detail, the packet's finite Euler-product inequality gives a positive lower bound $2^{-71}Q^{-7}$ for $|\zeta(\sigma+it)|$ when $|t|\ge1/64$, $Q\ge\log(|t|+3)$ and $1-2^{-88}Q^{-9}\le\sigma\le2$. The calculation uses $3+4\cos u+\cos(2u)=2(1+\cos u)^2$. Euler summation controls three derivatives and isolates the pole at one. A finite rectangular contour for the once-integrated prime-power sum, followed by monotone differences and finite Abel summation, yields the imported estimate
\begin{equation}\label{eq:growth-strong}
 n\ge 2^{\,2^{241}(m+5)^{10}}
 \quad\Longrightarrow\quad
 |\pi(n)-\mathrm{Li}_2(n)|\le 2^{-m} n/\log n .
\end{equation}
The exponent $10$ is strong enough to dominate every fixed power of $1/\log n$. The source gives rational approximations and precision bounds for the contour. We give the remaining deductions here.

\begin{lemma}[Finite factorial expansion]\label{lem:growthfactorial}
Let $r\ge1$, $C=2^{240}$, $M(r)=1+6400C^2r^2+48C$, $B_r=1+r!(1+2^{r+2})$, and
\[
 T_r=2^{\max(2M(r)^{10},\,32(r+2)^2)}.
\]
For $n\ge T_r$, set $L=\log n$, $u=1/L$, $P_r(u)=\sum_{j=0}^{r-1}j!u^j$. Then
\[
 \left|\frac{\pi(n)}{nu}-P_r(u)\right|\le B_r u^r.
\]
\end{lemma}
\begin{proof}
Write $D=2^{241}$, $\ell=\lfloor\log_2 n\rfloor$, and let $a$ be the largest integer with $D a^{10}\le\ell$, a bounded search. At $n\ge2^{D6^{10}}$ equation \eqref{eq:growth-strong} with $m=a-5$ applies. Since $\ell\ge L$ here and $a>(\ell/D)^{1/10}-1$, this bounds the normalized error $L|\pi-\mathrm{Li}_2|/n$ by $64\exp(-L^{1/10}/(2D))$, using $\log2\ge1/2$. Write $v=L^{1/10}$. For $n\ge2^{2M(r)^{10}}$ one has $v\ge M(r)$, so $\log64\le v/(4D)$ and $10r\log v\le10r\sqrt v\le v/(4D)$. Thus $|\pi-\mathrm{Li}_2|\le n/L^{r+1}$. The elementary inequality $\log v\le\sqrt v$ for $v\ge1$ follows by maximizing $2\log t-t$ at $t=2$; the signs of its derivative equally give the comparison on rational subintervals.

Telescoping the derivatives of $t/(\log t)^j$ yields
\[
 \mathrm{Li}_2(n)=nuP_r(u)
 -2\sum_{j=0}^{r-1}\frac{j!}{(\log2)^{j+1}}
 +r!\int_2^n\frac{dt}{(\log t)^{r+1}} .
\]
The finite boundary sum is at most $2^{r+1}r!$ in absolute value: bound $j!\le(r-1)!$ and sum a geometric progression (take $r=1$ separately). The integral split at $\sqrt n$ is at most $2^{r+1}\sqrt n+2^{r+1}n/L^{r+1}$. For $L\ge16(r+2)^2$, one has $(r+1)(\log2+\log L)\le L/2$, using $\log L\le\sqrt L$. Thus $2^{r+1}\sqrt n\le n/L^{r+1}$ and $n/L^{r+1}\ge1$. The three error terms give the stated $B_r$. Each integration identity follows at precision $s$ by telescoping on a uniform rational mesh. For $f(t)=t/(\log t)^j$, $|f''(t)|\le 2^{j+3}(j+1)^2$ on $t\ge2$, directly by differentiation. Choosing mesh width at most $2^{-s}/((n+1)2^{j+4}(j+1)^2)$ bounds the accumulated Taylor error by $2^{-s}$. The iteration in $r$ uses a finite list of these estimates with $s$ enlarged by the binary length of $2(r+1)!$.
\end{proof}

The same bound, with $B'_r=2^{r+1}B_r+\sum_{j<r}j!(j+2)$, holds for real $x$ supplied with its integer part whenever $x\ge2T_r$. Indeed $n=\lfloor x\rfloor\ge\sqrt{x}$ and $\pi(x)=\pi(n)$. For $t\in[n,x]$ the derivative of $t/(\log t)^{j+1}$ has magnitude at most $j+2$. Also $x/(\log x)^{r+1}\ge1$. Summing these bounds proves the assertion. The expansion for $\mathrm{li}(x)$ follows by the same integration identity with the additional constant $\mathrm{li}(2)$, bounded in absolute value by 4. The following eventual inequalities all have thresholds computed from these formulas.

\begin{proposition}[Reciprocal expansion and inequalities]\label{prop:growthrecip}
Define signed integers $k_m$ recursively, for $m\ge1$, by
\[
 k_m=m\,m!-\sum_{i=1}^{m-1}(m-i)!k_i.
\]
They start $1,3,13,71,461,3441$. For each $m$, with $u=1/\log x$,
\[
 \frac{x}{\pi(x)}
 =u^{-1}-1-\sum_{i=1}^m k_i u^i+O_m(u^{m+1}).
\]
The constants and thresholds in this expansion are primitive recursive in $m$.
\end{proposition}
\begin{proof}
Put $Q_m=1-u-\sum_{i=1}^m k_i u^{i+1}$ and $r=m+2$. In $P_r Q_m$ the linear coefficient vanishes. The coefficient at $u^{j+1}$, $1\le j\le m$, is $(j+1)!-j!-\sum_{i=1}^j(j-i)!k_i=0$ by definition. This explains the varying power $u^i$ in the formula. For a polynomial write $\|P\|_1$ for the sum of the absolute values of its coefficients. For $0<u\le1$ the residual is at most $(1+\|P_r\|_1\|Q_m\|_1)u^r$. Put $f=\pi(x)/(xu)$. Beyond $2T_r$ the preceding lemma now gives
\[
 |1-fQ_m|\le [1+(\|P_r\|_1+B'_r)\|Q_m\|_1]u^r.
\]
If also $x\ge2^{4(B'_r+1)}$, then $B'_ru^r\le1/2$ and $f\ge1/2$. Divide by $fu$ to obtain the claimed bound with twice the bracketed integer as constant.
\end{proof}

\paragraph{Consequences for inequalities on the source pages.}
For sufficiently large $x$, writing $L=\log x$, the following strict bounds hold:
\begin{align*}
 &\frac{x}{L}<\pi(x)<1.25506\frac{x}{L},\qquad
 \frac{x}{L-1}<\pi(x)<\frac{x}{L-1.1},\\
 &\frac{x}{L}\left(1+\frac1L+\frac2{L^2}\right)<\pi(x)
 <\frac{x}{L}\left(1+\frac1L+\frac2{L^2}+\frac{7.59}{L^3}\right),\\
 &\frac{x}{\pi(x)}< L-1-L^{-1}-3L^{-2}.
\end{align*}
For the first and third comparisons use the polynomial $1+u+2u^2+6u^3$ with error $B'_4u^4$. For the second, the reciprocal expansion starts $L-1-u+O(u^2)$, which lies strictly between $L-1.1$ and $L-1$. For the last comparison its next coefficient is $-13u^3$. In each use of an error $Cu^{j+1}$ against a signed leading term $au^j$ it suffices to add the integer threshold $x\ge2^{\lceil4(C+1)/|a|\rceil}$. These proofs establish eventual versions. The much smaller published numerical cutoffs require separate analysis or verified finite computations. The upper coefficient $1.25506$ alone also follows eventually from the basic PNT by taking $k=4$.

For integers $n$ the harmonic approximation in MathWorld is asymptotically equivalent to $\pi(n)$, since $0\le H_n-\log n\le1$ and the basic PNT applies. It eventually underestimates $\pi(n)$ strictly. Indeed, $H_n-\log(n+1)$ increases by $1/(n+1)-\log(1+1/(n+1))\ge0$. At $n=8$, $H_8>2.717$ while $\log9<2.2$ (the Taylor polynomial through degree five already gives $e^{1.1}>3$). Thus $H_n-3/2>L-1$ for $n\ge8$ and the lower bound $\pi(n)>n/(L-1)$ proves the assertion once $L>1$.

\subsection{Ramanujan's eventual inequality}
\begin{theorem}\label{thm:growthram}
For all sufficiently large positive rational $x$, and in the located-floor interpretation for real arguments,
\[
 \pi(x)^2<\frac{ex}{\log x}\pi(x/e).
\]
The same eventual inequality holds with $\mathrm{li}$ in both positions.
\end{theorem}
\begin{proof}
Put $u=1/\log x$, $P=1+u+2u^2+6u^3+24u^4$, $B=B'_5+4$, and take $x\ge6T_5$, $u\le1/2$. Both arguments lie in the range of the expansion. Divide both sides by $x^2u^2$. The left side is $P^2$ with error at most $(68+B)Bu^5$, since $\|P\|_1=34$. The right side is
\[
 \sum_{j=0}^4\frac{j!u^j}{(1-u)^{j+1}}+E,
 \qquad |E|\le64Bu^5.
\]
Indeed $\log(x/e)=1/u-1$. At each degree $d\le4$, the coefficient on the right is $\sum_{j=0}^d\binom{d}{j}j!$. The left coefficients are $1,2,5,16,64$ and the right coefficients are $1,2,5,16,65$. The discarded coefficients of $P^2$ sum to at most $34^2$. For each $j$, the geometric derivative series for $(1-u)^{-j-1}$ after degree $4-j$ has tail at most $64u^{5-j}$: factor out $u^{5-j}$ and compare with $2^{5-j}\sum_{h\ge0}\binom{h+j}{j}2^{-h}=64$. This identity also follows by multiplying $j+1$ geometric series with explicit tails. Thus the normalized right side minus the left is $u^4$ with error at most $Ku^5$, where $K=(68+B)B+64B+34^2+64\cdot34$. For $x\ge2^{4(K+1)}$ the difference is at least $u^4/2$. The same five coefficients and the chosen $B$ cover $\mathrm{li}$, since a constant shift of at most four is absorbed by the unnormalized remainder.

There is an explicit method to determine the floor of $x/e$ for $x=a/b\ge2$ with $a,b$ positive integers. Set $E_a=\sum_{i=0}^a1/i!$ and scan $1\le j\le\lfloor a/b\rfloor$. For these $j$ the number $a/(bj)$ is a multiple of $1/a!$. Since $0<e-E_a<1/a!$ for $a\ge2$ by the factorial tail, we have $j\le x/e$ exactly when $bjE_a<a$. Include $j=0$ and take the largest passing index. More precisely the factorial tail satisfies $1/(a+1)!<e-E_a<2/(a+1)!$. Since distinct points of the $1/a!$ mesh are $(a+1)/(a+1)!$ apart and $a\ge2$, this gives $|a-bje|\ge bj/(a+1)!$ for every tested $j$. The unnormalized gap established above is at least $x^2/(2(\log x)^6)\ge1/(2a^4)$. Thus a finite strict comparison is available: evaluate $e$ and $\log x$ with interval widths at most $1/[100(a+1)^6]$. On $L\ge2$ the reciprocal changes by at most the width of its argument, and $x\pi(x/e)\le a^2$. The width of the resulting interval for the difference is at most $5a^2/[100(a+1)^6]$, strictly smaller than the gap. The series for $e$ has remainder less than $1/j!$ after degree $j\ge2$, so this precision too has an explicit finite cutoff.
\end{proof}

\subsection{Inversion at the \(m\)-th prime}
\begin{theorem}[The first terms of the inverse expansion]
For $t=\log m$, $w=\log t$, the expansion on the Wikipedia page holds through both displayed correction terms:
\[
 \frac{p_m}{m}
 =t+w-1+\frac{w-2}{t}
 -\frac{w^2-6w+11}{2t^2}
 +O\!\left(\frac{w^3}{t^3}\right).
\]
All uses of the error below have primitive-recursive thresholds. In particular $p_m\sim m\log m$.
\end{theorem}
\begin{proof}
Write $x=p_m$, $R=x/m$, and $L=\log x=t+\log R$. Since $x\ge m$, the hypotheses of Proposition~\ref{prop:growthrecip} apply as soon as $m$ exceeds its threshold. At order two let $C_2$ be the explicit error constant obtained in its proof, so that
\[
 |R-G(R)|\le C_2/t^3,\qquad
 G(y)=t+\log y-1-\frac1{t+\log y}-\frac3{(t+\log y)^2}.
\]
Take $m$ large enough that $w\ge4$, $t\ge8w$ and $C_2/t^3\le1$. For instance $m\ge2^{2\max(256,C_2+1)}$ suffices for these size conditions, using $\log t\le\sqrt t$. Here $R\ge1$. The preceding estimate gives $L/2\le R\le2L$, hence $R\ge t/2$. If $R\ge4t$, then $\log R\le R/4$ (at $R\ge16$), so $R\le2t+R/2$ and $R\le4t$. Thus $R$ always lies in $[t/2,4t]$.

Set $a=w-1$, $b=w-2$, $c=-(w^2-6w+11)/2$, $Y=t+a+b/t+c/t^2$. Since $|a|,|b|\le w$ and $|c|\le2w^2$, the relative correction $d=Y/t-1$ satisfies $|d|\le2w/t\le1/4$. Expand $\log(1+d)$ through degree two; the geometric remainder is at most $2|d|^3$. This gives
\[
 \log Y=w+a/t+(b-a^2/2)/t^2+E_1,\qquad |E_1|\le30w^3/t^3 .
\]
Indeed $|d-a/t|\le2w/t^2$, bounding the change in the quadratic term by $3w^2/t^3$, and $|c|/t^3\le2w^2/t^3$. Also $z=\log Y$ satisfies $0\le z\le2w$ and $|z-w|\le3w/t$ by the same series. The identity
$(t+z)^{-1}=t^{-1}-z/t^2+z^2/(t^2(t+z))$ gives
\[
 \left|(t+z)^{-1}-t^{-1}+w/t^2\right|\le(3w+4w^2)/t^3,\qquad
 \left|(t+z)^{-2}-t^{-2}\right|\le4w/t^3.
\]
For the latter inequality use $t^{-2}-(t+z)^{-2}\le2z/t^3$. Substitution shows that $G(Y)$ has successive coefficients $w-1$, $a-1$, $b-a^2/2+w-3$, exactly $a,b,c$. The errors just bounded give $|G(Y)-Y|\le50w^3/t^3$. On the interval between $R$ and $Y$ the derivative
\[
 G'(y)=\frac1y\left(1+(t+\log y)^{-2}+6(t+\log y)^{-3}\right)
\]
is bounded by $1/2$ since $y\ge t/2$ and $t\ge8$. Integrating this explicit bound (or summing it on rational meshes) gives $|R-Y|\le |R-G(R)|+|G(R)-G(Y)|+|G(Y)-Y|\le(C_2+50)w^3/t^3+|R-Y|/2$. The asserted remainder bound is therefore $2(C_2+50)w^3/t^3$.
\end{proof}

This also gives the eventual versions of the displayed Rosser, Dusart and Axler bounds for $p_m$. To see the signs explicitly, write
\[
 F(g)=t+w-1+(w-2)/t-g/(2t^2).
\]
The expansion implies
\[
 F(w^2-6w+11.321)<p_m/m<F(w^2-6w+10.667).
\]
Indeed, with $K=2(C_2+50)$, it suffices to have $Kw^3/t<0.16$. One explicit extra condition is $t\ge(10^5(K+1))^2$: the inequality $w\le8t^{1/8}$ follows from $e^{w/8}\ge w/8$, and so $Kw^3/t\le512K/t^{5/8}<0.16$. For $w\ge12$, these bounds also imply the versions with $F(w^2)$ and $F(w^2-6w)$, the latter on the upper side. They give the upper bound $t+w-1+(w-2)/t$ because $w^2-6w+10.667>0$. They give the lower bound $t+w-1+(w-2.1)/t$ when $w^2-6w+11.321<0.2t$, which follows, for example, from $w^2\le64t^{1/4}$ at $t\ge 2^{20}$. Finally $w-2.1>0$ and $(w-2)/t<0.0513$ eventually (use $w\le\sqrt t$). Thus the familiar lower bounds $t$ and $t+w-1$, and upper bounds $t+w$, $t+w-0.5$ and $t+w-0.9487$, all follow with room to spare. A required real threshold $t\ge B$ is guaranteed by the integer test $m\ge2^{\lceil2B\rceil}$; $w\ge12$ can be ensured by $t\ge2^{24}$. For the two sharpest decimal comparisons the residual positive margin after division by $m$ is at least $(0.321/2-0.16)/t^2=1/(2000t^2)$. It is enough to approximate both $t,w$ within $1/[10^6(m+1)^4]$: in the region $t\ge2$, $0\le w\le m+1$, the sum of the absolute partial derivatives of $F(w^2-6w+c)$ with $10\le c\le12$ is less than $100(m+1)^2$ by differentiating the displayed rational expression. The logarithm enclosures above supply these widths, using $|(\log t)'|\le1$ for $t\ge1$ and one extra precision bit to propagate the first enclosure. Thus the strict conclusion can be verified by a prescribed positive separation in the object language. The small numerical onset of each individual inequality is an additional quantitative problem.

\section{The exact maximum of the normalized prime count}
\label{sec:chebsharp}
The numerical table up to a million has a consequence for all larger inputs. A factorial ratio bounds $\vartheta$ on the tail strongly enough to show that the maximum has already been reached.

\begin{theorem}
For rational $x>1$, and equally for real arguments with their integer parts supplied,
\[
 \frac{\pi(x)\log x}{x}\le \frac{30\log113}{113}<1.25506 .
\]
Equality in the first inequality occurs precisely at $x=113$ (in particular at rational cuts). Thus the upper Rosser--Schoenfeld inequality in the two source pages holds on its full stated range.
\end{theorem}
\begin{proof}
Put
\[
 g(u)=\lfloor u\rfloor-\lfloor u/2\rfloor-\lfloor u/3\rfloor
        -\lfloor u/5\rfloor+\lfloor u/30\rfloor .
\]
The increments of these five floors under translation by $30$ sum to zero. Each floor is determined by $\lfloor u\rfloor$, and for the residues $0$ through $29$ the successive values, in groups of five, are
\[
 01111\quad10111\quad01011\quad00101\quad00010\quad00001 .
\]
Thus $g$ is nonnegative and equals one on $[1,6)$. Define $F(n)=n!\lfloor n/30\rfloor!/(\lfloor n/2\rfloor!\lfloor n/3\rfloor!\lfloor n/5\rfloor!)$. Counting multiples of $p^j$ in each factorial gives
$v_p(F(n))=\sum_{j\ge1}g(n/p^j)$. This proves integrality and shows that every prime power $n/6<p^j\le n$ contributes a unit to this valuation. Taking logarithms yields
\[
 \psi(n)-\psi(\lfloor n/6\rfloor)\le\log F(n).
\]

For each $y>0$ let $B(y)=y\log y-y$. If $1\le y\le n$ and $m=\lfloor y\rfloor$, the lower and upper rectangles of $\log t$ between $1$ and $m$ give
$B(m)+1\le\log(m!)\le B(m)+1+\log m$. Furthermore $0\le B(y)-B(m)\le\log y$. When $0<y<1$ use $m=0$ and $-1\le B(y)\le0$ (integrate $\log t$, or use $y\log(1/y)\le1-y$). Hence in all cases
$|\log(\lfloor y\rfloor!)-B(y)|\le1+\log n$, for integers $n\ge1$. In the signed combination of the five $B(n/d)$ the terms involving $n\log n$ and $n$ cancel by $1-1/2-1/3-1/5+1/30=0$. The remaining coefficient of $n$ is
\[
 a=\tfrac12\log2+\tfrac13\log3+\tfrac15\log5-\tfrac1{30}\log30<0.93 .
\]
For instance $\log2<0.694$, $\log3<1.099$, $\log5<1.610$, and $\log30>3.400$ follow from the logarithm enclosures used for the finite table. The coefficient is positive, also directly from the formula and $30=2\cdot3\cdot5$. We obtain $\log F(n)\le a n+5(1+\log n)$. Repeated division by six takes at most $1+\log n/\log6$ nonzero stages; summing their arguments costs at most $6n/5$. Consequently
\[
 \psi(n)\le\tfrac65 a n+5(1+\log n)(1+\log n/\log6).
\]
With $b=10^6$, the second term divided by $n$ is bounded by $5(1+\log b)^2/b<1125/10^6$ for $n\ge b$. Indeed $\log6>1$, $\log b<14$, and $(1+\log t)^2/t$ decreases when $\log t>1$ by its derivative $(1+\log t)(1-\log t)/t^2$. This proves $\vartheta(t)\le\psi(t)<1.14t$ for every $t\ge b$ by taking the integer part.

The finite computation in \texttt{finite\_bounds.py} gives $\pi(b)=78498$, $\vartheta(b)>0.99b$, $\log b<13.82$ and $30\log113/113>1.25$. Subtract the partial-summation identity for $\pi(b)$ from that for $\pi(x)$, $x\ge b$. It gives
\[
 \pi(x)\le U(x):=\pi(b)-\frac{\vartheta(b)}{\log b}
       +\frac{1.14x}{\log x}
       +1.14\int_b^x\frac{dt}{\log^2 t}.
\]
At $x=b$, the normalized right side is strictly below $78498(13.82)/10^6-0.99+1.14<1.24$. Throughout the tail,
\[
 U'(x)=\frac{1.14}{\log x}
  < \frac{1.25}{\log x}\left(1-\frac1{\log x}\right)
   =\frac{d}{dx}\frac{1.25x}{\log x}.
\]
Here $\log b>12$ since $b>2^{19}$ and $\log2\ge2/3$ by the first term of the expansion of $2\int_0^{1/3}(1-u^2)^{-1}du$. Integration gives $\pi(x)<1.25 x/\log x$ on the entire tail. The derivative comparisons have the same finite meaning as the partial summations: on each compact interval cut a rational mesh, telescope increments of products and reciprocal logarithms, then use the rectangular enclosures for $1/t$ and $1/(t\log^2 t)$. All denominators on $[b,x]$ have the displayed positive lower bounds, and the initial strict margin at $b$ is independent of the mesh precision.

For integers $2\le n\le b$ the companion computes both endpoints of an enclosure of $\log n$ and compares $113\pi(n)\log n$ with $30n\log113$. Every case except $n=113$ is strictly smaller; there $\pi(113)=30$. The same computation gives $30\log113/113<1.25506$. For $1<x<2$ the count vanishes, and on $[2,3)$ we have $\log x/x<1$. On $[3,\infty)$ the multiplier $\log x/x$ decreases strictly. Indeed for $v>u\ge3$, $\log v-\log u\le(v-u)/u<(v-u)\log u/u$. Between integer jumps the largest value is therefore at the left endpoint. For a named real with known integer part $113$, equality in the normalized bound forces $x=113$: a positive separation from $113$ would give the strict decrement just proved, and the lower bound $x\ge113$ makes the absence of positive separation equivalent to equality of its real name. This finishes both the inequality and its equality statement.
\end{proof}

\section{The ordinary M\"obius series for Riemann's smooth count}
\label{sec:gram}

The derivative in MathWorld (11) contains a conditionally convergent constant term. There is also a normalization constant in Wikipedia's formula for $R$. We evaluate both in the ordinary order $n=1,2,\ldots$. The one analytic number-theory input in this section is the quantitative $\vartheta$ estimate in the ProveIt manuscript used in Section~\ref{sec:growthrefine}. Its status as an imported human proof applies to these deductions.

\begin{lemma}[Two boundary values of the M\"obius series]\label{lem:Rboundary}
Write $M(x)=\sum_{n\le x}\mu(n)$, with $n\ge1$ in all sums. With primitive-recursive moduli of convergence,
\[
 \sum_{n\ge1}\frac{\mu(n)}n=0,\qquad
 \sum_{n\ge1}\frac{\mu(n)\log n}{n}=-1 .
\]
\end{lemma}
\begin{proof}
\emph{A bound for the prime-power sum.}
Write $H(x)=1+\log x$ for $x\ge1$. We first obtain
\[
 |\psi(n)-n|\le W n/H(n)^8,\qquad W=2^{3000}100!.
\]
The imported estimate from the packet's \texttt{theta\_bridge.tex} is
$|\vartheta(n)-n|\le2^{-r}n$ for $n\ge2^{C(r+1)^{10}}$, $C=2^{240}$.
Let $\ell=\lfloor\log_2 n\rfloor$. If $\ell\ge C$, take the largest $r$ with $C(r+1)^{10}\le\ell$. Then $H(n)\le2C(r+2)^{10}$. Since $e^z\ge z^{80}/80!$ and $\log2\ge1/2$, we have $2^{-r}(r+2)^{80}\le 3\cdot 2^{80}80!$. If $\ell<C$, use $H(n)\le C+2$ and $\vartheta(n)\le\psi(n)\le 4n$. This last bound follows from $L_n\le2^{4n}$ in Section~\ref{sec:dyadicgrowth}, since $\psi(n)=\log L_n$. Finally $\psi(n)-\vartheta(n)\le\sqrt n\log n$, since each prime $p\le\sqrt n$ contributes at most $\log n$ in its higher powers. This gap times $H(n)^8/n$ is at most $2^{10}9!$ by applying the exponential bound to $H(n)/2$. The stated $W$ covers both contributions.

\emph{Passing from $\psi$ to $M$.}
For $i=0,\ldots,8$ put $m=i/2$, $K_0=1$, $K_{i+1}=2^{50}(1+W+K_i)$. We prove $|M(x)|\le K_i x/H(x)^m$ at integer $x\ge1$. The case $i=0$ follows from $|\mu|\le1$. Put $E(u)=\psi(u)-u$ for integer $u$. The two exact convolutions
\[
 \mu(n)\log n=-\sum_{d\mid n}\mu(d)\Lambda(n/d),
 \qquad \sum_{d\le x}\mu(d)\lfloor x/d\rfloor=1
\]
give
\begin{equation}
 M(x)\log x
 =-1-\sum_{d\le x}\mu(d) E(\lfloor x/d\rfloor)
       +\int_1^x M(t)\frac{dt}{t}.
 \label{eq:gram-m}
\end{equation}
For the first convolution, a squarefree $n$ contributes one term for each prime divisor, with $\mu(n/p)=-\mu(n)$. When precisely one prime $p$ has exponent $v\ge2$, its surviving terms $n/d=p^v,p^{v-1}$ cancel. Two repeated primes make every term zero. The second convolution counts the divisors of each integer up to $x$ with weight $\mu$. The integral expands as the finite sum $\sum_{n\le x}\mu(n)\log(x/n)$.

For $\ell=\lfloor\log_2 x\rfloor\ge64$ set
$Y=2^{\lfloor\sqrt\ell\rfloor}$ and $D=\lfloor x/(Y+1)\rfloor$. Write $H=H(x)$. The elementary comparisons
\[
 D\ge\sqrt x,\qquad \sqrt H/4\le\log(Y+1),\qquad
 1+\log Y\le3\sqrt H
\]
follow from $\ell/2\le\log x\le\ell+1$ and $\sqrt\ell\le\ell/4$. For instance $Y\le x^{1/4}$ and $x/(x^{1/4}+1)\ge\sqrt x+1$ on this range. In the part $d\le D$ the argument of $E$ is at least $Y+1$. Summing the bound $W(x/d)4^8/H^4$ costs at most $4^8 W x/H^3$, using $\sum_{d\le x}1/d\le H$.

For $d>D$, expand $E(q)=\sum_{u\le q}(\Lambda(u)-1)$ and interchange the finite sums. The resulting expression is
\[
 \sum_{u\le Y}(\Lambda(u)-1)
                   (M(\lfloor x/u\rfloor)-M(D)).
\]
Both arguments of $M$ are at least $D$ and at most $x/u$. The inductive bound makes the difference at most $2^{m+1}K_i x/(uH^m)$ in absolute value. Partial summation and $\psi(u)\le4u$ give
$\sum_{u\le Y}(\Lambda(u)+1)/u\le9(1+\log Y)$.
Thus this part costs at most $864K_i xH^{1/2-m}$ for $m\le4$.
In the integral in \eqref{eq:gram-m}, the portion up to $\sqrt x$ costs at most $\sqrt x$. Beyond $\sqrt x$ we have $H(\lfloor t\rfloor)\ge H/3$, so the remainder is at most $3^m K_i x/H^m$. Finally $\sqrt x\le2^6 4!\,x/H^4$ by the exponential bound.
Dividing by $\log x\ge H/2$, these estimates give $K_{i+1}x/H^{m+1/2}$ for $i\le7$. For $\ell<64$, $H^4\le65^4<2^{50}$ suffices. Applying this improvement eight times gives the asserted bound. Each application proves its prefix estimates at a fixed ceiling and precision. There is no object induction on a formula quantifying over every ceiling.

\emph{Ordinary convergence and evaluation.}
Put $K=K_8$. On $t\ge3$, $|M(t)|\le16Kt/H(t)^4$ also for nonintegral $t$, by passing to $\lfloor t\rfloor$. For integers $V>N\ge3$ apply partial summation on $(N,V]$ to the weights $1/t$ and $(\log t)/t$. The two endpoint terms together are bounded by $2K/(\log N)^3$ for the second weight, and the integral by
\[
 16K\int_N^V\frac{dt}{t(\log t)^3}\le 8K/(\log N)^2 .
\]
The first weight has smaller endpoint bounds and integral at this range. Thus $32K/(\log N)^2$ bounds each signed tail. In particular $N_j=2^{2^{j+8}(K+1)}$ suffices for precision $2^{-j}$. Partial summation here means summation by parts on integer intervals followed by the integrals of the explicit weights. Their derivatives are rational combinations of $1/t$ and $\log t$, so these finite identities equally follow by rectangular sums as in Section~\ref{sec:weights}.

For rational $0<s\le1$, define $Z=\sum_{n\ge1}n^{-1-s}$, $Z_1=\sum_{n\ge1}(\log n)n^{-1-s}$. Absolute convergence and the divisor cancellation $\sum_{d\mid n}\mu(d)=[n=1]$ give
\[
 \sum_n\frac{\mu(n)}{n^{1+s}}=\frac1Z,\qquad
 \sum_n\frac{\mu(n)\log n}{n^{1+s}}=-\frac{Z_1}{Z^2}.
\]
For the second equality multiply the identity $\sum_{d\mid n}\mu(d)=[n=1]$ by $\log n=\log d+\log(n/d)$, then sum with weight $n^{-1-s}$. Absolute convergence follows with effective tails: if $s=a/b$, blocks $2^{bj}\le n<2^{b(j+1)}$ have total weight at most $b2^{-j}$ without a logarithm and $b^2(j+1)2^{-j}$ with one. Cutting a product at $mn\le T^2$ misses only pairs having a coordinate $>T$, so these estimates justify the convolution products as well.

Comparison with the integral yields $1\le sZ\le1+s$. Also $|s^2 Z_1-1|\le2s^2$. Indeed $f(t)=t^{-1-s}\log t$ has one maximum, at $\log t=1/(1+s)$, of height at most one. The difference between its integral and unit-width rectangular sum is at most its total variation, at most two. Integration of $f$ gives $1/s^2$, either by integration by parts or by the substitution $t=e^u$. The tail in this calculation is the endpoint expression $T^{-s}(1+s\log T)/s^2$. It follows that $1/Z\le s$ and $|Z_1/Z^2-1|\le5s$.

Lastly the ordinary limits equal the damped limits as $s$ decreases to zero. On the first $N$ terms the discrepancy caused by $n^{-s}$ is bounded by $sN^3$ for either weight, using $1-e^{-v}\le v$. On a tail whose every prefix sum is bounded by $\epsilon$, multiplying by the decreasing sequence $n^{-s}\in[0,1]$ leaves its sum bounded by $\epsilon$, by finite summation by parts. First choose $N=N_j$, then a rational $s\le 2^{-j}/N^3$. This proves both values and gives a primitive-recursive comparison throughout.
\end{proof}

\begin{theorem}[Gram series, M\"obius formula and derivative]
For $x>1$ at rational input, let $t=\log x$. The following series define the same real name. The first extends to $x=1$ with value one.
\[
 R(x)=1+\sum_{j\ge1}\frac{t^j}{j\,j!\,\zeta(j+1)}
      =\sum_{n\ge1}\frac{\mu(n)}n\operatorname{li}(x^{1/n}).
\]
Furthermore,
\[
 R'(x)=\frac1{x\log x}\sum_{n\ge1}\frac{\mu(n)}n x^{1/n}.
\]
\end{theorem}
\begin{proof}
For $u>0$, write
$\operatorname{li}(e^u)=\gamma+\log u+\sum_{j\ge1}u^j/(j\,j!)$,
with the principal-value normalization. The exponential series integrated after extracting its constant term gives this expansion (the constant cancels below). At $u=t/n$, its logarithmic term equals $\log t-\log n$. Lemma~\ref{lem:Rboundary} annihilates $\gamma+\log t$ and evaluates the remaining logarithmic sum as $+1$. The double series of positive powers is absolutely convergent, since its absolute contribution for a given $n$ is at most $(e^t-1)/n^2$. Divisor cancellation at the absolutely convergent exponent $j+1\ge2$ gives
$\sum_n\mu(n)n^{-j-1}=1/\zeta(j+1)$. This proves the first display.

Differentiating the Gram polynomials in $t$ gives
$\sum_{j\ge1}t^{j-1}/(j!\zeta(j+1))$. Multiplying by $t$ and expanding the reciprocal zeta factors gives $\sum_n \mu(n)(e^{t/n}-1)/n$. The same $n^{-2}$ bound permits this interchange. Adding back $\sum\mu(n)/n=0$ and using $dt/dx=1/x$ proves the derivative formula, with the ordinary ordering of $n$.

For quantitative details, on $|t|\le A$ with integer $A\ge1$, the Gram series of polynomials and its first two derivatives have tails bounded by the corresponding exponential tails (for derivative order $d$, truncate the exponential after degree $J-d$). For $m\ge2A$ the successive majorant terms $A^m/m!$ decrease by a factor at least two; their sum through any index is at most $e^A\le3^A$. Taking $J=5A+j+8$ therefore gives all three tails below $2^{-j}$. A Taylor estimate for the finite polynomials bounds linearization error by $3^A h^2$ when both arguments lie in the chosen interval. These estimates prove the stated derivative as a modulus assertion on rational increments. For any rational $x>1$ a working bound is $A=\lceil x\rceil+2$ on a sufficiently small neighbourhood. The absolute $n$-tails above are at most $3^A/N$ after $N$ terms, and the two boundary tails have the explicit bound in Lemma~\ref{lem:Rboundary}. All limits and rearrangements thus reduce to bounded operations on finite rational enclosures.
\end{proof}

\section{The exact cutoff for the logarithmic-integral inequality}
\label{sec:li-cut}
MathWorld gives Ramanujan's inequality for $\pi$ and a companion inequality for the principal-value logarithmic integral. The companion has an exact integer transition. Monotonicity reduces the whole infinite range to two rational sign certificates.

\begin{theorem}
For integers $n\ge2$,
\[
 \operatorname{li}(n)^2 <
 \frac{en}{\log n}\operatorname{li}(n/e)
 \quad\Longleftrightarrow\quad n\ge2418 .
\]
The inequality is reversed strictly for $2\le n\le2417$.
\end{theorem}
\begin{proof}
For $x>e$ put $l=\log x$, $f(x)=\operatorname{li}(x)$, and
\[
 D(x)=e f(x/e)-\frac{l}{x} f(x)^2 .
\]
Multiplying $D(x)$ by $x/l>0$ gives the right side of the claimed inequality minus its left side. The derivative $f'(x)=1/l$ gives
\[
 D'(x)=\frac1{l-1}-\frac{2f(x)}x+(l-1)\frac{f(x)^2}{x^2}
       =\frac{(1-(l-1)f(x)/x)^2}{l-1}\ge0.
\]
Consequently it suffices to determine the signs at $2417$ and $2418$ on the interval $[3,\infty)$. For $n=2$, $0<2/e<1$ and $f(2/e)<0$ by integration of a strictly negative integrand. Strictness of the two endpoint certificates below suffices even if the square in the derivative vanishes somewhere.

Here are the definitions and numerical error bounds behind the sign test. With $t=\log x$ write
\[
 \operatorname{Ei}(t)=\operatorname{li}(x)
   =\operatorname{PV}\int_{-\infty}^{t} e^u\,du/u .
\]
Substitution $x=e^u$ gives the displayed normalization. Symmetric excisions at $x=1$ or at $u=0$ have the same limit because $\log(1+h)/[-\log(1-h)]\to1$; this follows by sandwiching each log between $h/(1+h)$ and $h/(1-h)$. On finite intervals the Taylor series of $(e^u-1)/u$ can be integrated termwise with the following explicit bound. Write $S_K(t)=\sum_{j=1}^{K}t^j/(j\,j!)$. For $0<t\le B=32$ and $K=256$,
\[
 \operatorname{Ei}(t)=\log(t/B)+S_K(t)-S_K(-B)+E,\qquad
 -\frac{2^{-B}}{B}-\epsilon\le E\le\epsilon,\quad
 \epsilon=\frac{4B^{K+1}}{(K+1)!}.
\]
Indeed the integral to $-B$ is negative and has size at most $e^{-B}/B$. The principal value of $1/u$ from $-B$ to $t$ is $\log(t/B)$. Each series tail after integration is bounded by $2B^{K+1}/(K+1)!$ since successive absolute terms have ratio at most $1/2$. These arguments also provide arbitrary-precision names by increasing $B$ and $K$.

To evaluate logarithms normalize a positive rational to $2^a y$ with $1\le y\le2$. With $z=(y-1)/(y+1)$ integrate the finite geometric series for $2/(1-z^2)$. This gives
\[
 0\le \log y-2\sum_{j=0}^{49}\frac{z^{2j+1}}{2j+1}
 \le \frac{2z^{101}}{101(1-z^2)}.
\]
The same rule evaluates $\log 2$. Evaluate $e$ by summing through $1/50!$ with tail at most $2/51!$. For $t=\log n$, the quantity $nD(n)=en\operatorname{Ei}(t-1)-t\operatorname{Ei}(t)^2$ has the following enclosures. All endpoints in the table are integers divided by $1000$.
\[
\begin{array}{c|rr}
 n & \text{strict lower numerator}& \text{strict upper numerator}\\\hline
 2417 & -437 &-435\\
 2418 & 40 &42
\end{array}
\]
The file \texttt{certify\_li.py} reproduces these rational comparisons. It accumulates $t^j/j!$ by multiplication by $t/j$ and rounds the two endpoints out to multiples of $2^{-120}$. It applies the logarithm bounds at each end and uses monotonicity of $\operatorname{Ei}$ on the positive axis. Thus each line is a closed finite calculation, reproducible using integer arithmetic.

For clarity, the monotonicity step also has a finite interpretation. Fix integers $3\le a\le b$. On $[3,b]$, $\log x-1\ge1/16$; already the positive terms $2(1/2+(1/2)^3/3)$ of the logarithm series for $3$ exceed $1+1/16$. We have $|\operatorname{li}(x)|\le b+10$ there. Indeed the same expansion at $t=\log3$ gives $|\operatorname{li}(3)|<10$ (also checked by the auxiliary script), and $|f'|\le1$ on $[3,b]$. Write $A=1-(l-1)f/x$. The elementary derivative bounds give $|A|\le 1+b(b+10)$ and $|A'|\le (b+1)(b+11)$ using $x\ge3$, $l\le b$. Differentiating $A^2/(l-1)$ thus gives $|D''|\le M=10^4(b+12)^5$. On a mesh of width $h$ the error of the left sum of $D'$ for $D(b)-D(a)$ is at most $M(b-a)h$. This estimate can also be derived without a mean-value witness: differences of the integral names $f$ and $\log$ on $[u,u+h]$ equal their left derivatives times $h$ with quadratic remainders, by upper and lower rectangles for their decreasing positive kernels. Applying the product and reciprocal difference rules gives the same first-order expansion for $D$; equivalently integrate the bounded derivative of $D'$ cell by cell. Take $h\le2^{-s-2}/(M(b+1))$ and give each summand an accumulated rounding allowance below $2^{-s-2}$. The sum of squares is nonnegative, so $D(b)\ge D(a)-2^{-s}$ at every requested precision. Fixed positive margins in the two endpoint tests turn this into the stated strict inequalities. The mesh, all apartness bounds, and the power-series lengths are primitive recursive in the endpoints and precision.
\end{proof}

For an explicit general series schedule in the preceding paragraph, let the required absolute error be $2^{-s}$, take $B\ge s+\lceil|t|\rceil+8$, and use $K+1=2m$ with $m=4(B+1)^2+s+8$. Since $(2m)!\ge m^m$ and $m\ge4B^2$, the tail allowance $4 B^{2m}/(2m)!$ is at most $4^{1-m}$. The logarithm at $t>0$ in these applications is evaluated with the known lower bound on $t$ (at $x/e$ it is at least $1/16$ for $x\ge3$). General dyadic widths can be propagated by the finite circuit estimate of Section~\ref{sec:mellin}.

\section{Exact numerical consequences and two source corrections}
\label{sec:cert}
\subsection{Fourteen decades}
For $x=10^e$ the last three columns below give nearest integers, with every rounding decided by an interval lying strictly within a single rounding cell. The first difference is signed as in Wikipedia, the second as $\operatorname{li}(x)-\pi(x)$, and the last as in MathWorld.
\begin{center}
\begin{tabular}{r r r r r}
\toprule
$e$ & $\pi(x)$ & $\pi(x)-x/\log x$ & $\operatorname{li}(x)-\pi(x)$ & $R(x)-\pi(x)$\\
\midrule
$1$ & $4$ & $0$ & $2$ & $1$ \\
$2$ & $25$ & $3$ & $5$ & $1$ \\
$3$ & $168$ & $23$ & $10$ & $0$ \\
$4$ & $1\,229$ & $143$ & $17$ & $-2$ \\
$5$ & $9\,592$ & $906$ & $38$ & $-5$ \\
$6$ & $78\,498$ & $6\,116$ & $130$ & $29$ \\
$7$ & $664\,579$ & $44\,158$ & $339$ & $88$ \\
$8$ & $5\,761\,455$ & $332\,774$ & $754$ & $97$ \\
$9$ & $50\,847\,534$ & $2\,592\,592$ & $1\,701$ & $-79$ \\
$10$ & $455\,052\,511$ & $20\,758\,029$ & $3\,104$ & $-1\,828$ \\
$11$ & $4\,118\,054\,813$ & $169\,923\,159$ & $11\,588$ & $-2\,318$ \\
$12$ & $37\,607\,912\,018$ & $1\,416\,705\,193$ & $38\,263$ & $-1\,476$ \\
$13$ & $346\,065\,536\,839$ & $11\,992\,858\,452$ & $108\,971$ & $-5\,773$ \\
$14$ & $3\,204\,941\,750\,802$ & $102\,838\,308\,636$ & $314\,890$ & $-19\,200$ \\
\bottomrule
\end{tabular}
\end{center}
The count column is an evaluation of the fourth-root formula in Section~\ref{sec:survivors}, implemented in \texttt{prime\_table.cpp}. Each approximate square- or cube-root seed is corrected until two integer quotient inequalities certify the root. An independent computation verifies the first ten rows as follows. Store $S(v)=v-1$ at the distinct positive quotients $v=\lfloor x/j\rfloor$. At step $p$, $S(v)$ counts the primes already processed and the integers in $[2,v]$ that survived them. Process $p$ when $S(p)>S(p-1)$, up to $\lfloor\sqrt{x}\rfloor$. The new removals for $v\ge p^2$ are $pm$ with $p\le m\le v/p$ surviving earlier primes. Subtract $S(\lfloor v/p\rfloor)-S(p-1)$ in descending $v$ so both counts still have their old values. All required quotients remain stored since iterated quotients compose. At termination every composite has been removed and $S(x)=\pi(x)$.

For the analytic columns use Section~\ref{sec:li-cut} and the Gram power series in Section~\ref{sec:gram}. In this table $\log x<34$. For the coefficient $\zeta(j+1)$, truncate after $u=68\cdot2^{\lceil14/j\rceil}$ terms and add the upper tail $1/(ju^j)$ from integrating $t^{-j-1}$. Rounding each positive summand down on a $2^{-150}$ grid introduces at most $u2^{-150}$ further uncertainty. Use the directed products for $(\log x)^j/j!$ through $j=256$ and majorize the remaining series by $2\cdot34^{257}/257!$. The exponential-integral computation uses the same Taylor remainder with $34$ in place of $32$ when necessary. The file \texttt{approximation\_table.py} checks every rounding against this table. Only rows through $10^{14}$ are claimed as computations here.

\subsection{Finite certificates for bounds and the harmonic approximation}
Write $L=\log n$, $h_n=H_n-3/2$ and $U_3(n)=n(L^{-1}+L^{-2}+2L^{-3})$. The sieve and the rational logarithm checks give the following universal statements about the indicated \emph{integer} intervals. Section~\ref{sec:chebsharp} supplies a global upper bound separately.
\begin{center}\small
\begin{tabular}{@{}p{.62\linewidth}l@{}}
\toprule
Verified comparison & Range inclusive\\ \midrule
$n/L<\pi(n)$ & $17\le n\le10^6$\\
$n/(L-1)<\pi(n)$ & $5393\le n\le10^6$\\
$\pi(n)<n/(L-1.1)$ & $60184\le n\le10^6$\\
$U_3(n)\le\pi(n)$ & $88789\le n\le10^6$\\
$\pi(n)\le U_3(n)+7.59n/L^4$ & $2\le n\le10^6$\\
$n/h_n<\pi(n)$ & $1429\le n\le10^6$\\
\bottomrule
\end{tabular}
\end{center}

\textbf{Counterexamples at the published lower endpoint.}
The displayed three-term lower inequality on the Wikipedia page starts at $88783$. The sieve gives $\pi(n)=8596$ for $88783\le n\le88788$. The enclosure at the first integer gives
\[
 8596.04<U_3(88783)<8596.05 .
\]
All six inputs violate the inequality, with a positive certified margin in each case. A change from a strict to a non-strict inequality would still give six failures. The interval in the table begins with the next prime. The eventual bound is proved in Section~\ref{sec:growthrefine}.

Here is how to reproduce the arithmetic in these comparisons. For $n\ge1$ write $n=2^a y$, $1\le y<2$; take $j=\lfloor1024y\rfloor$, $y=(j/1024)(1+z)$. Precompute the $1025$ grid logarithms by the rational series of Section~\ref{sec:li-cut}. Bound the small remainder by $z-z^2/2\le\log(1+z)\le z$ by integrating $1-t\le1/(1+t)\le1$. Add $a\log2$ and round outward to a $2^{-52}$ grid. Multiply each inequality by its positive denominator, using upper and lower endpoints in opposite directions to certify the sign. For example, if $l,u$ enclose $L$, a sufficient check for $U_3(n)\le\pi(n)$ is $\pi(n)l^3\ge n(u^2+u+2)$. To refute it use the strict reverse with the endpoints exchanged. For harmonic sums the prefix of $\sum_{i\le n}\lfloor2^{52}/i\rfloor\,2^{-52}$ is a lower bound. The code \texttt{finite\_bounds.py} gives every loop limit and integer comparison. Its inductive invariant is the correct prime prefix, the directed harmonic and logarithmic bounds, and a bit recording that every comparison so far has passed.

MathWorld describes the harmonic approximation on $50\le n\le1000$ as accurate to about two. An exact rational scan shows how to read that description. The minimum signed error $\pi(n)-n/h_n$ is at $n=222$ and lies in $(-2.530,-2.529)$. Its maximum is at $n=887$ and lies in $(2.779,2.780)$. For instance these results follow by maintaining the exact pair for $H_n=H_{n-1}+1/n$ and comparing cross products at each step. The first seven positive errors occur at $1,109,113,114,199,200,201$, with $n=2$ omitted since $h_2=0$. For Ramanujan's \emph{prime-count} inequality the successes among integers $2,\ldots,18$ are $6,9,10,12,14,15,16,18$. The checker encloses $e$ between the Taylor polynomial through degree 12 and that polynomial plus $2/13!$, decides the floors of $n/e$, and cross-multiplies using logarithm enclosures. The exact all-integer cutoff $2418$ proved earlier belongs to the logarithmic-integral inequality.

\paragraph{Finite checks for the \(k\)-th prime.}
The ten displayed lower or upper bounds for $p_k$ whose stated starting index is at most $10^6$ were also tested at every integer from that index through $10^6$. There is exactly one exception. The Wikipedia page uses $0.9487$ at $k=39017$, but
\[
 p_{39017}=467473,\qquad
 467471<39017(\log39017+\log\log39017-0.9487)<467472.
\]
Dusart's original coefficient is $0.9484$; this passes the same entire interval. For Axler's $10.667$ upper bound, the separate indices $46254381$, $46254382$, and $46255381$ pass. The first follows, for example, since the proposed upper bound exceeds $905060142$ and $\pi(905060142)=46254381$. These counts use an Eratosthenes sieve and, for the three large endpoints, the quotient sieve. Nested logarithms use monotonicity and outward-rounded rational intervals. For the quadratic term, $(w-3)^2+c-9$ gives endpoint bounds even on intervals crossing $w=3$. The exact tests and their outputs are in \texttt{checks/nth\_bounds.py} and \texttt{verification/nth\_bounds.txt}. The original constant appears in \href{https://doi.org/10.1090/S0025-5718-99-01037-6}{P.~Dusart, \emph{The $k$th prime is greater than $k(\ln k+\ln\ln k-1)$ for $k\ge2$}}.

\section{Two conditional prime-gap implications}
\label{sec:conditional-gap}
Let $c=\pi_{\rm c}$ denote the circle constant and let $x$ be an integer. A positive difference $\pi(x)-\pi(y)$ supplies a prime by scanning the integers in $(y,x]$. The question is which explicit finite inequality suffices to make this difference positive.

\paragraph{Two values of the counting error.}
Suppose $2657\le y<x$ are integers and at these two arguments the Schoenfeld bounds quoted on the Wikipedia page hold,
\[
 |\pi(t)-\operatorname{li}(t)|<\frac{\sqrt t\log t}{8c},\qquad t=y,x .
\]
If $x-y\ge\sqrt x\,(\log x)^2/(4c)$, then $(y,x]$ contains a prime. Indeed, $\operatorname{li}(x)-\operatorname{li}(y)\ge(x-y)/\log x$ by a lower rectangle, while the two endpoint errors in total are strictly smaller than $\sqrt{x}\log x/(4c)$. This gives a pointwise implication with no zero-counting premise.

A shorter interval requires information on how the error varies across it. To see the scale, one can construct an integer counting function with jumps $0,1$ that obeys the same error envelope and has an arbitrarily far-out gap of length at least $2\sqrt{x}\log x$. Start at $m_0=2657$ with a rational approximation to $\operatorname{li}(m_0)$ of error at most $1/4$. At every integer $m\ge m_0$ add a nonnegative rational lower approximation to $1/\log m$ of error at most $2^{-m}$; then take the integer part of each accumulated rational. This gives $G(m)$ with unit jumps and $|G(m)-\operatorname{li}(m)|<4$, since the difference between upper and lower unit rectangles telescopes to at most $1/\log m_0$. The envelope exceeds four throughout: $c\le4$, $\sqrt{2657}>50$ and $\log2657>11/2$. Now fix $x\ge2^{4096}$, put $S=\sqrt x$, $L=\log x$, and choose an integer $H$ between $2SL$ and $2SL+2$ by taking the ceiling of a rational upper enclosure of width at most one. Put $y=x-H$. Freeze $G$ at $G(y)$ through $x$ and subtract $G(x)-G(y)$ from every later value. Since $H\le3SL$ and $L\le 24x^{1/24}$ (integrate $1/t\le t^{1/24-1}$), $y\ge x/2$ and $\log y\ge L/2$. The added discrepancy after $y$ is bounded by $8+H/\log y\le 8+6S$. The envelope there is at least $SL/128$, larger than the total error $12+6S$ since $L\ge2048$. Thus endpoint subtraction alone cannot yield a general statement about jumps at the shorter scale.

\paragraph{A finite local hypothesis yielding Dudek's constant.}
There is a more informative arithmetic premise that does give $4/c$ with one logarithm. Define a triangular sum, with $1\le h<t$ integers,
\[
 W(t,h)=\sum_{|m-t|<h}\bigl(1-|m-t|/h\bigr)\Lambda(m).
\]
It can also be computed as $[P(t+h)-2P(t)+P(t-h)]/h$, where
$P(u)=\sum_{m\le u}(u-m)\Lambda(m)$. Expand the coefficient of each $\Lambda(m)$ on the three intervals to get the identity. The following local inequality is the specific input isolated from \href{https://arxiv.org/abs/1402.6417}{Dudek, \emph{On the Riemann hypothesis and the difference between primes}, Section 2.2}:
\begin{equation}\label{eq:local-triangle}
 W(t,h)>h-\frac{2}{c}\left(\frac{t}{\sqrt{t-h}}
           +\frac{(t+h)^{3/2}}{t}\right)\log(2t/h)-\frac{48}{5h}.
\end{equation}
Dudek derives this premise from estimates for zeta zeros under RH when $2t/h>15$. We prove the following implication entirely with bounded arithmetic and real names.

\begin{proposition}
For each integer $x\ge2^{4096}$ choose $h$ with $(2/c)\sqrt x\log x-2\le h\le(2/c)\sqrt x\log x$ and put $t=x-h$. If \eqref{eq:local-triangle} holds at this pair, a bounded search returns a prime in
\[
 (x-(4/c)\sqrt x\log x,\ x].
\]
\end{proposition}
\begin{proof}
Write $d=2/c$, so $1/2\le d\le1$. Compute $h$ by taking the floor of a rational lower enclosure to $dSL$ of width at most one. Put $z=h/x\le L/S$. The bound $L\le24x^{1/24}$ gives $z\le1/8$ and $3zL\le1$ in the stated range; also $h\ge dSL/2$ and $2t/h>15$ (indeed $z<1/9$). Hence
\[
 \frac{t}{\sqrt{t-h}}+\frac{(t+h)^{3/2}}t\le S(2+6z),
 \qquad \log(2t/h)\le L/2-\log L+3 .
\]
For the first inequality use $(1-2z)^{-1/2}\le(1-2z)^{-1}\le1+4z$ and $(1-z)^{-1}\le1+2z$. For the second use $2t/h\le(4/d)S/L$ and $\log 8\le3$. Since $L\ge2048$ and $\log L\ge11/2$, the product on the right is at most $S(L-2\log L+6+3zL)$. Substitution in \eqref{eq:local-triangle} gives
\[
 W(t,h)>dS(2\log L-7)-12\ge 2S-12>S.
\]
We bound all proper prime powers in this window. The number of squares in $(x-2h,x]$ is at most
$1+2h/(\sqrt x+\sqrt{x-2h})\le1+2L$. For each $a\ge3$ there are at most $x^{1/3}$ eligible bases, and at most $L/\log 2\le2L$ exponents. Every weight is at most $L$. Thus the total proper-power contribution divided by $S$ is at most
\[
 (1+2L)L/S+2L^2 x^{-1/6}<1.
\]
For the strict bound insert $L\le24 x^{1/24}$; the three resulting terms are at most $24x^{-11/24}$, $1152x^{-5/12}$ and $1152x^{-1/12}$ respectively, whose sum is below one for $x\ge2^{4096}$. If all prime indicators in $(x-2h,x)$ were zero, the finite weight sum would therefore be $<S$. A bounded scan must succeed, and $2h\le(4/c)SL$ gives the claimed location. All comparisons with strict analytic hypotheses can be read as rational interval certificates with positive margin. This proves the arithmetic extraction from \eqref{eq:local-triangle}; deriving that premise from a formalized RH requires the analytic zero estimates.
\end{proof}

\section{From the arguments to Lean derivations}
\label{sec:lean}
Two theorem types play different roles here. A theorem about \texttt{Nat} in Lean checks a mathematical implementation in a familiar setting. A term of type $\Gamma\vdash[\mathsf{PRAAx}]A$ constructs a proof in the first-order object calculus. The input functions in a uniform object theorem are fixed PR definitions, or terms supplied by a metatheoretic schema; a general evaluator for all PR programs is unnecessary.

\subsection{The checked arithmetic components}
The archive has standalone Lean sources for trial division, rank selection, midpoint inversion, the factorial kernel and polynomial coefficients. They use Lean 4.19.0 and its standard library. In \texttt{Counting.lean}, \texttt{prime\_spec} identifies the Boolean test with the proper-divisor predicate. For any marked sequence, \texttt{select\_spec} proves that whenever the positive target rank has been reached by the supplied bound, the selected location is marked, has that exact rank, and has predecessor rank one less. The induction has two cases: the old prefix already reaches the target, or the new bit must be true and raise the count by precisely one. \texttt{select\_unique} establishes uniqueness of that marked location. The object proof still needs the Euclidean bound in Section~\ref{sec:found} to turn this bounded specification into unrestricted inverse rank.

\texttt{Inversion.lean} implements divisor-incidence inversion for integer sums. It proves the divisor sum of the recursively defined Möbius function, and proves the scaled \texttt{midpoint\_inversion} identity, parametrized by a Boolean marking function. Section~\ref{sec:powers} explains the correspondence of the scaled rows to $J_{1/2}$. The marking may be instantiated with \texttt{PrimeCounting.prime}. \texttt{ChebyshevKernel.lean} proves periodicity, the bit range, and the interval where the five-floor kernel equals one. \texttt{Growth.lean} verifies the two finite polynomial coefficient lists that decide the eventual sign of Ramanujan's inequality. \texttt{Sieve.lean} proves inclusion--exclusion for arbitrary lists of Boolean marks. Its \texttt{floor\_sieve} gives the alternating floor sum with subset lcm for arbitrary lists of moduli, including repetitions. For distinct primes the lcm equals the product by factorization. These files establish their indicated statements directly. They do not mechanize the analytic limits. See \texttt{verification/lean.txt} for the build and axiom output.

\subsection{The pinned object language}
Use ProveIt revision \texttt{a9c0f1949be61679ed2e219b97dc42ecb7ae91a9}, whose root project specifies Lean 4.32.0. The optional \texttt{ProveItPNT.lean} is a source-level client of the actual object theorem in \texttt{FinalPNT.lean}. Its type is the sequent of Section~\ref{sec:pntgrowth}. The included standalone verification used 4.19.0, so this client and the 4.32.0 ProveIt dependency closure have not been rebuilt here. The stronger logarithmic-integral error estimate and its $\vartheta$ version are reused from the separate human proof packet. Its authors describe its Lean material as an outline. Conclusions depending on these quantitative estimates carry that dependency even when their finite coefficient calculations have been checked.

In the following construction, \texttt{PTm n} is an arithmetic term in a context of $n$ number variables, and \texttt{Fm PTm n} is a formula. Compile every bounded scan into a primitive recursion storing its loop state and fixed parameters; project the needed component at the end. To use induction on a relation with a positive rational denominator, multiply both sides by the recorded denominators. Represent signed quantities by two nonnegative parts, and express $a-b\le c-d$ as $a+d\le c+b$. A finite interval or product can have its own denominator at each step; cross multiplication preserves order. The prime-power identity can instead use the common denominator $2(L!)^2$. Boolean universal and existence tests are justified by the prefix-product and first-hit invariants of Section~\ref{sec:found}.

The entries below give the proof obligations in dependency order. Module names in the last column are existing providers or statements to implement as indicated.
\begin{enumerate}
\item \textbf{Prime data.} Use \texttt{Valuation.isP} and the shipped theorems \texttt{isP\_of\_prime}, \texttt{prime\_of\_isP}, \texttt{isP\_01}. Use \texttt{PiTheta.Pi x zero} for the inclusive sum. Compile $E(k)$ and the bounded least scan. Express its minimality by the product bit for earlier failed comparisons. Obtain a successful prime search from the factorial argument before applying the rank lemma. Use \texttt{FactorTrace.factor\_product} or \texttt{FactorVector.reconstruct} for the positive factorization, and \texttt{factor\_injective} for reconstruction from equal valuation vectors.
\item \textbf{Survivor counts.} Define $\Phi$ directly by a double bounded sum/product. Define the subset mask product $d_s$ and parity fold. Prove the mask expansion by prefix induction in the integer-pair ring. Induct on $\lfloor N/d\rfloor$ and the final remainder to count the multiples. For $P_k$, compile the least-factor trace, its non-unit length and minimum index. Implement the two maps in \eqref{eq:sv-factor} by quotient and multiplication. To count carried indices, first sum their graph bits by source, then by target (finite double-sum swap). This yields the two explicit counts and both large sieve formulas. The wheel bound and first-crossing scan then need only signed inequalities. In the factorial congruence use the bounded inverse search modulo $p$ and pair off inverse residues.
\item \textbf{Weights and inversion.} Define the cut bits by $Bp^j<A$ and $Bp^j=A$; use the same $A,B$ throughout the $j=dk$ substitutions. Fill the Möbius table by increasing $n$, or use the positive and negative parts in the repository. Induct on $n$ to identify the divisor-recurrence version with the squarefree version, using their common divisor sum and isolating $d=n$. Existing helpers include \texttt{DivisorSwap.div\_swap}, \texttt{MobiusInversion.mobius\_inversion}, \texttt{PrimePowerMangoldt.lam\_pw} and \texttt{ChebyshevAbel.abel\_gen}. Transport the finite scaled inversion by clearing its denominator and swapping the two bounded sums. Form the lcm and prime product as integer codes before applying a log name.
\item \textbf{Analytic names.} Implement signed rational intervals for the four arithmetic operations. Inversion requires a named rational lower bound on the absolute denominator. Define $\log$ by rectangles or the range-reduced atanh series, $\exp$ by its factorial series, real powers by positive roots or $\exp(s\log x)$, and complex numbers by pairs of such enclosures. Define $\operatorname{li}$ by the finite exponential-integral expression in Section~\ref{sec:li-cut} with its tail. Define $\zeta$, $R$ and the step integrals by the cutoffs in Sections~\ref{sec:mellin}, \ref{sec:trivial} and \ref{sec:gram}. For each displayed equality, compile the finite expressions at requested precision $j+g$ on both sides and prove the width/error inequality for precision $j$. At a bounded fold fix the final rounding budget first, split it over all iterations, and induct on a numerical error accumulator. The circuit sensitivity bound and the explicit tail bounds supply $g$. This avoids induction on equality of real names, which would quantify over all precisions.
\item \textbf{Bounds and asymptotics.} Encode $\log t$, $\vartheta(t)$ and $\psi(t)$ by intervals with common requested accuracy. For the sharp global upper bound, prove the integer periodic kernel first, then the logarithm-of-factorial rectangle inequalities. Algebraically telescope the bounds for $\vartheta$ and use the finite data at $10^6$. For the PNT instantiate \texttt{arith\_pnt} and prove the change from harmonic mesh to the log name by the one-cell error bound. For higher estimates, the remaining substantial object-development dependency is translation of the stronger manuscript's finite Euler-product, contour and desmoothing estimates. Given those, the factorial polynomials, residual coefficient norms, eight finite improvements for $M$, and inverse substitutions in this report use the stated, bounded schedules. For $x/e$ use the factorial comparison in Section~\ref{sec:growthrefine}.
\item \textbf{Certificate reflection.} Compile the finite range checks as one state recursion: at each input update the prime count, lower/upper sums and a bit for all earlier comparisons. Prove once that its update preserves the intended arithmetic specification. For a fixed final height a numeral derivation reducing the bit to one furnishes a closed certificate. Large inputs are better replayed in blocks: store boundary states, verify each block and compose the state equalities by transitivity. The supplied Python/C++ runs evaluate the arithmetic; generating these object-calculus numeral derivations is further work.
\end{enumerate}

\paragraph{Proof scope.}
The full Riemann explicit formulas over the nontrivial zeros, their recommended summation orders and any interchange with the Möbius sum need analytic information beyond the Mellin identity in its half-plane of absolute convergence. This report proves the trivial-zero correction separately. Littlewood-type sign changes, zero-free regions at the sharper exponents quoted on Wikipedia, record-size table entries above $10^{14}$, and the published time and storage exponents of optimized sieves are outside the proved results. Section~\ref{sec:conditional-gap} specifies exactly which error bounds imply its short-interval conclusions. For non-global numerical lower bounds and estimates on $p_m$, the finite windows and the eventual results leave an intervening range unless a separate argument covers it.

\section*{Sources and reproduction}
\addcontentsline{toc}{section}{Sources and reproduction}
\begin{itemize}
\item \href{https://en.wikipedia.org/wiki/Prime-counting_function}{Wikipedia, \emph{Prime-counting function}}, and \href{https://mathworld.wolfram.com/PrimeCountingFunction.html}{Eric W. Weisstein, \emph{Prime Counting Function}, MathWorld}. The two requested entries identify the target formulas, endpoint conventions, numerical claims and tables, accessed November 22, 2026.
\item \href{https://github.com/VladimirReshetnikov/ProveIt/tree/a9c0f1949be61679ed2e219b97dc42ecb7ae91a9/Logic/PrimitiveRecursiveArithmetic}{Vladimir Reshetnikov, \emph{ProveIt}, pinned iPRA tree}. Relevant providers are the ConciseIPRA proof calculus, the ordinary derivations in \texttt{AnalyticNumberTheory}, and the \href{https://github.com/VladimirReshetnikov/ProveIt/tree/a9c0f1949be61679ed2e219b97dc42ecb7ae91a9/Logic/PrimitiveRecursiveArithmetic/docs/proofs/analytic-number-theory/pnt_li_ipra_2026-11-15}{logarithmic-integral packet}. The text identifies the two imported quantitative inputs separately.
\item \href{https://doi.org/10.1215/ijm/1255631807}{J. B. Rosser and L. Schoenfeld, \emph{Approximate formulas for some functions of prime numbers}}, Illinois J. Math. 6 (1962), 64--94. \href{https://doi.org/10.1007/s11139-016-9839-4}{P. Dusart, \emph{Explicit estimates of some functions over primes}}, Ramanujan J. 45 (2018), 227--251. Section~\ref{sec:chebsharp} gives an elementary constructive proof of the specific global upper bound.
\item \href{https://www-users.cse.umn.edu/~odlyzko/doc/arch/meissel.lehmer.pdf}{J. C. Lagarias, V. S. Miller and A. M. Odlyzko, \emph{Computing $\pi(x)$: the Meissel--Lehmer method}}, Math. Comp. 44 (1985), 537--560. \href{https://www.jstor.org/stable/2312732}{S. W. Golomb, \emph{On the ratio of $N$ to $\pi(N)$}}, Amer. Math. Monthly 69 (1962), 36--37.
\item \href{https://arxiv.org/abs/1407.1901}{A. Dudek and D. Platt, \emph{On solving a curious inequality of Ramanujan}}, Exp. Math. 24 (2015), 289--294. \href{https://arxiv.org/abs/1402.6417}{A. Dudek, \emph{On the Riemann hypothesis and the difference between primes}}. The conditional local weighted estimate is identified at \eqref{eq:local-triangle}.
\end{itemize}

\end{document}
